\chapter{Homotopia, nakrycia i komórkowe modele przestrzeni}
\label{ch:topologia-algebraiczna-podstawy}
\index{grupa podstawowa}

W Rozdziale~\ref{ch:algebra-homologiczna} kompleks łańcuchowy
był ciągiem modułów i homomorfizmów. Aby zbudować taki
kompleks z przestrzeni, nauczymy się rozkładać przestrzeń
na proste części. Równolegle zapytamy, kiedy dwa odwzorowania
można ciągle deformować jedno w drugie. Są to pytania
topologiczne: gładkość nie jest tu potrzebna.

Zaczniemy od dróg i nakryć, następnie od sympleksów
przejdziemy do CW-kompleksów. Homologia tych obiektów
pojawi się w następnym rozdziale. Homotopii odwzorowań
nie należy mylić z homotopią \emph{łańcuchową}
z Definicji~\ref{def:homotopia-lancuchowa};
ta druga jest algebraicznym odpowiednikiem pierwszej.

\section{Homotopie, drogi i grupa podstawowa}
\label{sec:homotopie-pi1}

W całym rozdziale $I=[0,1]$. Parametr $t\in I$
można interpretować jako czas ciągłej deformacji.
Obrazy pośrednie mogą być bardzo różne, lecz
ciągłość wyklucza skoki i rozdarcia.

\begin{definicja}[Homotopia i równoważność homotopijna]
\index{homotopia i rownowaznosc homotopijna@Homotopia i równoważność homotopijna}
\label{def:homotopia-topologiczna}
Dwa ciągłe odwzorowania $f,g:X\to Y$ są
\emph{homotopijne}, zapisujemy $f\simeq g$,
jeśli istnieje ciągłe $H:X\times I\to Y$ z
$H(x,0)=f(x)$ i $H(x,1)=g(x)$.
Jeśli $f|_A=g|_A$ oraz $H(a,t)=f(a)$
dla $a\in A\subseteq X$ i $t\in I$,
mówimy o homotopii \emph{względem $A$}.
Przestrzenie $X,Y$ są \emph{homotopijnie
równoważne}, gdy istnieją ciągłe $f:X\to Y$,
$g:Y\to X$ z $gf\simeq\id_X$ i
$fg\simeq\id_Y$.
\end{definicja}

Homeomorfizm ma prawdziwą odwrotność.
Równoważność homotopijna dopuszcza odwrotność
z dokładnością do deformacji. Dysk i punkt
nie są homeomorficzne, ale są homotopijnie
równoważne.
Na przykład dla $D^n=\{x\in\R^n:\|x\|\leq1\}$ odwzorowanie
$r:D^n\to\{0\}$ jest stałe, inkluzja
$i:\{0\}\hookrightarrow D^n$ spełnia $ri=\id_{\{0\}}$,
a $H(x,t)=(1-t)x$ łączy $\id_{D^n}$ z $ir$.
Homotopia odbywa się wewnątrz dysku, gdyż
$\|H(x,t)\|=(1-t)\|x\|\leq1$.

\begin{definicja}[Retrakt deformacyjny]
\index{retrakt deformacyjny@Retrakt deformacyjny}
\label{def:retrakt-deformacyjny}
Podprzestrzeń $A\subseteq X$ jest
\emph{silnym retraktem deformacyjnym} $X$,
jeśli istnieje ciągłe $H:X\times I\to X$ takie, że
 $H(x,0)=x$, $H(x,1)\in A$ dla $x\in X$
 oraz $H(a,t)=a$ dla $a\in A$, $t\in I$.
Odwzorowanie $r(x):=H(x,1)$ jest wówczas
retrakcją $X\to A$. Oznaczając inkluzję przez
$i:A\hookrightarrow X$, mamy $ri=\id_A$ i $ir\simeq\id_X$;
stąd $i$ i $r$ są odwrotnościami homotopijnymi.
Przestrzeń homotopijnie równoważną punktowi
nazywamy \emph{ściągalną}. Równoważnie, niepusta
przestrzeń $X$ jest ściągalna wtedy i tylko wtedy,
gdy $\id_X$ jest homotopijna z pewnym odwzorowaniem
stałym $X\to X$: odwrotność homotopijna odwzorowania
$X\to\{*\}$ wybiera właśnie punkt, do którego
ściągamy $X$.
\end{definicja}

\begin{przyklad}[Promieniowa deformacja]
\label{prz:retrakcja-pierscien}
Niech $n\geq1$. Dla $x\in\R^n\setminus\{0\}$ wzór
$H(x,t)=(1-t)x+t\,x/\|x\|$ daje silny retrakt deformacyjny
$\R^n\setminus\{0\}$ na $S^{n-1}$.
Istotnie, $H(x,t)=((1-t)+t/\|x\|)x$, a współczynnik
jest dodatni, więc tor nie przechodzi przez usunięty
początek. Dla $\|x\|=1$ punkt pozostaje nieruchomy,
a dla $t=1$ otrzymujemy $x/\|x\|$. W płaszczyźnie
punkty spoza okręgu poruszają się do środka, zaś
punkty między okręgiem a usuniętym początkiem
poruszają się na zewnątrz, zawsze wzdłuż promienia.
Każdy niepusty zbiór wypukły $C\subseteq\R^N$ jest ściągalny:
dla wybranego $c_0\in C$ stosujemy
$H(x,t)=(1-t)x+tc_0$, którego cały obraz należy
do $C$ z definicji wypukłości.
\end{przyklad}

\begin{figure}[!htbp]
\centering
\begin{tikzpicture}[>=Stealth,line cap=round,line join=round,scale=.95]
 \fill[manifoldblue!5] (-2.7,-1.75) rectangle (2.7,1.75);
 \draw[manifoldblue,line width=1.7pt] (0,0) circle(1.25);
 \fill[white] (0,0) circle(3pt);
 \draw[manifoldred,line width=1pt] (0,0) circle(3pt);
 \node[manifoldred,below right=2pt] at (0,0) {$0$};
 \fill[manifoldred] (1.95,1.13) circle(2.5pt);
 \draw[manifoldred,very thick,-{Stealth[length=2.8mm]}]
   (1.95,1.13)--(1.08,.62);
 \fill[manifoldgold!90!black] (-.55,-.32) circle(2.5pt);
 \draw[manifoldgold!90!black,very thick,-{Stealth[length=2.8mm]}]
   (-.55,-.32)--(-1.08,-.62);
 \node[manifoldblue!90!black,above right=2pt]
   at (.70,1.04) {$S^1$};
 \node[manifoldred,right=3pt] at (1.95,1.13) {$x$};
 \node[manifoldgold!80!black,above left=2pt]
   at (-.55,-.32) {$y$};
\end{tikzpicture}
\caption{Promieniowa retrakcja $\R^2\setminus\{0\}$ na $S^1$.
Punkty $x$ i $y$ przesuwają się w przeciwne strony
promieni, a żaden tor nie przechodzi przez $0$.}
\label{fig:retrakcja-promieniowa}
\end{figure}

Do badania dziur potrzebujemy dróg,
których końce pozostają nieruchome w czasie deformacji.

\begin{definicja}[Drogi, pętle i homotopia dróg]
\index{drogi petle i homotopia drog@Drogi, pętle i homotopia dróg}
\label{def:droga-petla}
\emph{Droga} w $X$ to ciągłe
$\alpha:I\to X$; jej początek i koniec
to $\alpha(0)$ i $\alpha(1)$.
Gdy oba są równe $x_0$, nazywamy ją
\emph{pętlą opartą w $x_0$}.
Drogi $\alpha,\beta$ o tych samych końcach
są homotopijne \emph{względem końców},
jeśli łączy je homotopia $H:I\times I\to X$
z $H(0,t)=\alpha(0)$ i
$H(1,t)=\alpha(1)$ dla każdego $t$.
Gdy $\alpha(1)=\beta(0)$, określamy
\[
 (\alpha*\beta)(s)=
 \begin{cases}
  \alpha(2s),&0\leq s\leq\tfrac12,\\
  \beta(2s-1),&\tfrac12\leq s\leq1,
 \end{cases}
\]
Przy $s=1/2$ oba wzory dają ten sam punkt,
więc lemat o sklejaniu zapewnia ciągłość.
Ścieżka odwrotna to $\bar\alpha(s)=\alpha(1-s)$.
Symbol $\alpha*\beta$ oznacza: najpierw $\alpha$,
potem $\beta$; parametr czasu każdej z dróg
przebiega dwa razy szybciej.
\end{definicja}

Homotopia względem końców jest relacją równoważności.
Stała deformacja daje zwrotność, zamiana $t$ na $1-t$ daje
symetrię, a sklejenie dwóch homotopii na połówkach przedziału
czasu daje przechodniość; ciągłość wynika z lematu o sklejaniu.

\begin{przyklad}[Drogi w zbiorze wypukłym]
Jeżeli $C\subseteq\R^N$ jest wypukły, a drogi
$\alpha,\beta:I\to C$ mają te same końce,
to $H(s,t)=(1-t)\alpha(s)+t\beta(s)$ jest homotopią
względem końców. Wypukłość gwarantuje $H(s,t)\in C$,
zaś $H(0,t)=\alpha(0)$ i $H(1,t)=\alpha(1)$.
Na przykład dwie krzywe w dysku prowadzące od
$(-1,0)$ do $(1,0)$, choćby przez przeciwną połowę
dysku, można połączyć taką deformacją.
Warunek ustalenia końców jest potrzebny przy
składaniu klas dróg: bez niego środkowy punkt
$\alpha(1)=\beta(0)$ mógłby podczas deformacji
przestać być wspólnym końcem.
\end{przyklad}

\begin{figure}[!htbp]
\centering
\begin{tikzpicture}[>=Stealth,line cap=round,line join=round,scale=.9]
 \fill[manifoldblue!5,rounded corners=18pt]
  (-.4,-1.15) rectangle (10.2,2.45);
 \draw[manifoldblue!35,thick,rounded corners=18pt]
  (-.4,-1.15) rectangle (10.2,2.45);
 \coordinate (x) at (.5,.25);
 \coordinate (y) at (9.25,.25);
 \draw[manifoldred,very thick,-{Stealth[length=2.8mm]}]
  (x) .. controls (2.6,2.4) and (6.7,2.0) .. (y)
  node[pos=.48,above=4pt] {$\alpha$};
 \draw[manifoldgreen!80!black,very thick,-{Stealth[length=2.8mm]}]
  (x) .. controls (2.7,-1.0) and (6.5,-.7) .. (y)
  node[pos=.51,below=4pt] {$\beta$};
 \draw[manifoldgold!90!black,thick,-{Stealth[length=2.5mm]}]
  (4.9,1.5) -- (4.9,.35)
  node[pos=.52,right=4pt,font=\small] {deformacja};
 \fill[manifoldblue] (x) circle(2.4pt) (y) circle(2.4pt);
 \node[below left=3pt] at (x) {$x_0$};
 \node[below right=3pt] at (y) {$x_1$};
\end{tikzpicture}
\caption{Homotopia dróg przesuwa ich wnętrza,
ale zachowuje oba punkty końcowe.}
\label{fig:homotopia-drog}
\end{figure}

\begin{lemat}[Reparametryzacja nie zmienia klasy]
\label{lem:reparametr-drogi}
Jeśli $r:I\to I$ jest ciągła i spełnia
$r(0)=0$, $r(1)=1$, to
$\alpha\circ r$ jest homotopijna z $\alpha$
względem końców.
\end{lemat}
\begin{proof}
Połóżmy $r_t(s)=(1-t)r(s)+ts$.
Każde $r_t$ jest ciągłe, ma te same końce
i wartości w $I$. Funkcja
$H(s,t)=\alpha(r_t(s))$ jest ciągła.
Dla $t=0$ daje $\alpha\circ r$, dla $t=1$
daje $\alpha$, a końce pozostają stałe.
\end{proof}

Pętla jest zamkniętą drogą, a jej klasa pamięta,
czy można ją odkształcić do pętli stałej bez
przesuwania punktu bazowego. Składanie pętli
pozwala zamienić to pytanie geometryczne na
rachunek w grupie.

\begin{twierdzenie}[Grupa podstawowa]
\label{thm:grupa-podstawowa}
Klasy homotopii pętli opartych w $x_0$,
z działaniem $[\alpha][\beta]=[\alpha*\beta]$,
tworzą grupę $\pi_1(X,x_0)$.
Element neutralny to stała pętla $c_{x_0}$,
a odwrotnością $[\alpha]$ jest $[\bar\alpha]$.
\end{twierdzenie}
\begin{proof}
Jeśli $H$ deformuje $\alpha$ w $\alpha'$,
a $K$ deformuje $\beta$ w $\beta'$
względem końców, sklejenie
\[
 L(s,t)=
 \begin{cases}
   H(2s,t),&s\leq\tfrac12,\\
   K(2s-1,t),&s\geq\tfrac12
 \end{cases}
\]
jest ciągłe, bo na wspólnej krawędzi
obie części mają wartość $x_0$.
Działanie nie zależy więc od reprezentantów.

Przy trzech pętlach $(\alpha*\beta)*\gamma$
przeznacza na $\alpha,\beta,\gamma$
odcinki czasu długości $1/4,1/4,1/2$,
zaś $\alpha*(\beta*\gamma)$ długości
$1/2,1/4,1/4$. Obie drogi przebiegają
te same kawałki w tej samej kolejności
i różnią się zmianą parametru. Dokładniej,
$\alpha*(\beta*\gamma)=((\alpha*\beta)*\gamma)\circ r$,
gdzie $r(s)=s/2$ na $[0,1/2]$,
$r(s)=s-1/4$ na $[1/2,3/4]$ oraz
$r(s)=2s-1$ na $[3/4,1]$.
Wartości tych trzech wzorów zgadzają się
w punktach podziału, a $r(0)=0$, $r(1)=1$.
Lemat~\ref{lem:reparametr-drogi} daje łączność
na klasach. Podobnie $c_{x_0}*\alpha=\alpha\circ r_-$
dla $r_-(s)=\max(0,2s-1)$ oraz
$\alpha*c_{x_0}=\alpha\circ r_+$
dla $r_+(s)=\min(2s,1)$.
Obie reparametryzacje zachowują końce,
więc $[c_{x_0}]$ jest elementem neutralnym.

Dla odwrotności określmy
\[
 H(s,t)=
 \begin{cases}
  \alpha(2(1-t)s),&0\leq s\leq\tfrac12,\\
  \alpha(2(1-t)(1-s)),&\tfrac12\leq s\leq1.
 \end{cases}
\]
Przy $s=1/2$ obie gałęzie są równe
$\alpha(1-t)$, a przy $s=0,1$ mają
wartość $x_0$. Przy $t=0$ otrzymujemy
$\alpha*\bar\alpha$, przy $t=1$ pętlę stałą;
jest to zatem homotopia względem punktu bazowego.
Zamiana $\alpha$ z $\bar\alpha$ daje też
$[\bar\alpha][\alpha]=[c_{x_0}]$.
\end{proof}

\begin{stwierdzenie}[Odwzorowania indukowane i punkt bazowy]
\label{prop:pi1-funktor}
Odwzorowanie ciągłe $f:(X,x_0)\to(Y,y_0)$
z $f(x_0)=y_0$ wyznacza homomorfizm
\[
 f_*:\pi_1(X,x_0)\to\pi_1(Y,y_0),
 \qquad [\alpha]\mapsto[f\circ\alpha].
\]
Zachodzi $(gf)_*=g_*f_*$ i
$(\id_X)_*=\id$. Homotopia odwzorowań
utrzymująca $x_0$ w $y_0$ daje
równe homomorfizmy. Jeśli
$u$ jest drogą od $x_0$ do $x_1$,
to $b_u([\alpha])=[\bar u*\alpha*u]$
jest izomorfizmem
$\pi_1(X,x_0)\to\pi_1(X,x_1)$.
\end{stwierdzenie}
\begin{proof}
Złożenie $f\circ H$ przenosi homotopię
pętli w homotopię ich obrazów, a
$f\circ(\alpha*\beta)
=(f\circ\alpha)*(f\circ\beta)$.
Równości dla złożeń wynikają stąd wprost.
Jeśli $F$ jest homotopią odwzorowań utrzymującą
punkt bazowy, $F(\alpha(s),t)$ jest
homotopią pętli opartych w $y_0$.
Jeżeli $\alpha$ i $\alpha'$ są homotopijne
względem końców, dołączenie na początku
$\bar u$ i na końcu $u$ do wszystkich
pośrednich dróg daje homotopię pętli przy $x_1$.
Zatem $b_u$ nie zależy od reprezentanta.
Ponadto $b_u([\alpha])b_u([\beta])$ jest
reprezentowane przez
$\bar u*\alpha*u*\bar u*\beta*u$.
Środkowa droga $u*\bar u$ jest homotopijna
względem końców z drogą stałą przy $x_0$,
więc po usunięciu postoju otrzymujemy
$\bar u*(\alpha*\beta)*u$.
To dowodzi własności homomorfizmu.
Z tego samego rachunku
$b_{\bar u}b_u=\id_{\pi_1(X,x_0)}$
oraz $b_u b_{\bar u}=\id_{\pi_1(X,x_1)}$.
\end{proof}

Izomorfizm $b_u$ wymaga drogi łączącej oba punkty:
wszystkie pętle przy $x_0$ pozostają w jego
składowej łukowej. Wybór drogi może mieć znaczenie.
Dla innej drogi $u'$ od $x_0$ do $x_1$ połóżmy
$\lambda=\bar u*u'$, pętlę przy $x_1$.
Ponieważ $u*\lambda\simeq u'$ względem końców,
z definicji wynika
$b_{u'}([\alpha])=[\lambda]^{-1}b_u([\alpha])[\lambda]$.
Stąd przy zmianie drogi izomorfizm może zmienić
się o sprzężenie w $\pi_1(X,x_1)$.

\begin{wniosek}[Niezmienniczość grupy podstawowej]
\label{cor:pi1-homotopia}
Jeśli $f\simeq g:X\to Y$ przez homotopię
$H$, a $u(t)=H(x_0,t)$ prowadzi od
$f(x_0)$ do $g(x_0)$, to
$g_*=b_u\circ f_*$. Każda równoważność
homotopijna $f:X\to Y$ daje więc
izomorfizm grup podstawowych,
po uwzględnieniu zmiany punktu bazowego.
\end{wniosek}
\begin{proof}
Piszmy $H_t(x)=H(x,t)$ i dla $t\in I$ połóżmy
$u_t(s)=u(ts)$. Droga $u_t$ prowadzi od $f(x_0)$
do $H_t(x_0)$. Dla pętli $\alpha$ przy $x_0$
rodzina $(u_t*(H_t\circ\alpha))*\bar u_t$ jest
ciągłą homotopią pętli opartych w $f(x_0)$:
ciągłość na częściach przedziału wynika z ciągłości
$H$, a na granicach części wzory się zgadzają.
Przy $t=0$ otrzymujemy klasę $[f\circ\alpha]$,
gdyż $u_0$ jest stała. Przy $t=1$ otrzymujemy
$[u*(g\circ\alpha)*\bar u]
=b_{\bar u}([g\circ\alpha])$.
Stąd $f_*=b_{\bar u}g_*$, czyli $g_*=b_u f_*$.

Niech $h:Y\to X$ będzie odwrotnością homotopijną
$f$ i niech $y_0=f(x_0)$. Dla jasności piszemy
$f_{*,a}$, gdy $a$ jest punktem bazowym w dziedzinie.
Homotopia $hf\simeq\id_X$ wyznacza drogę
$u$ od $h(y_0)$ do $x_0$ oraz równość
$b_u\,h_{*,y_0}f_{*,x_0}=\id$.
Zatem $f_{*,x_0}$ jest iniektywny.
Podobnie homotopia $fh\simeq\id_Y$ wyznacza
drogę $v$ od $fh(y_0)$ do $y_0$ i równość
$b_v(fh)_{*,y_0}=\id$. Ponieważ $b_v$ jest
izomorfizmem, $(fh)_{*,y_0}$ jest surjektywny,
a więc także $f_{*,h(y_0)}$ jest surjektywny.
Obraz drogi $u$ przez $f$ prowadzi od $fh(y_0)$
do $y_0$. Dla pętli $\beta$ przy $h(y_0)$
tożsamość $f\circ(\bar u*\beta*u)
=\overline{f\circ u}*(f\circ\beta)*(f\circ u)$
daje naturalność zmiany punktu bazowego:
$f_{*,x_0}b_u=b_{f\circ u}f_{*,h(y_0)}$.
Prawa strona jest surjektywna, a $b_u$ jest
izomorfizmem; dlatego $f_{*,x_0}$ jest
surjektywny i w konsekwencji jest izomorfizmem.
\end{proof}

Silna retrakcja deformacyjna
z~\ref{def:retrakt-deformacyjny} daje
izomorfizm grup podstawowych w punktach
podprzestrzeni. Dla przestrzeni wypukłej
każdą pętlę $\alpha$ można deformować
wzorem $(1-t)\alpha(s)+tx_0$ do pętli
stałej, nie ruszając punktu bazowego.
Stąd $\pi_1(C,x_0)$ jest grupą trywialną
dla każdego niepustego wypukłego $C\subseteq\R^N$
i $x_0\in C$.

\begin{przyklad}[Pętla na brzegu dysku]
\label{prz:petla-w-dysku}
W dysku $D^2\subseteq\R^2$ weźmy punkt bazowy
$x_0=(1,0)$ i pętlę
$\alpha(s)=(\cos(2\pi s),\sin(2\pi s))$.
Wzór $F(s,t)=(1-t)\alpha(s)+t x_0$ daje
homotopię do pętli stałej: nierówność
$\|F(s,t)\|\leq(1-t)\|\alpha(s)\|+t\|x_0\|=1$
pokazuje, że obraz pozostaje w dysku, a
$F(0,t)=F(1,t)=x_0$ utrzymuje punkt bazowy.
W połowie deformacji punkt naprzeciw $x_0$
przechodzi przez środek, bo $F(1/2,1/2)=0$.
Zatem ta konkretna deformacja nie mieści się
w $D^2\setminus\{0\}$; w Sekcji~\ref{sec:nakrycia}
obliczymy, że na samym okręgu $S^1$ nie istnieje
żadna homotopia tej pętli do stałej.
\end{przyklad}

\section{Nakrycia: lokalnie wiele kopii jednej przestrzeni}
\label{sec:nakrycia}

Nakrycie przypomina nałożenie kilku arkuszy
na tę samą małą okolicę. Gdy idziemy po bazie,
wybranie początkowego arkusza jednoznacznie
ustala przebieg drogi na nakryciu.
Po zamkniętej drodze możemy wrócić na inny arkusz.

\begin{definicja}[Nakrycie i arkusze]
\index{nakrycie i arkusze@Nakrycie i arkusze}
\label{def:nakrycie}
Surjekcja ciągła $p:E\to X$ jest \emph{nakryciem},
jeśli każdy $x\in X$ ma otwarte otoczenie $U$
takie, że
\[
 p^{-1}(U)=\coprod_{\lambda\in\Lambda}V_\lambda,
 \qquad
 p|_{V_\lambda}:V_\lambda\xrightarrow{\cong}U,
\]
gdzie $V_\lambda$ są parami rozłącznymi
otwartymi podzbiorami $E$, a ograniczenia $p|_{V_\lambda}$
są homeomorfizmami.
Takie $U$ nazywamy \emph{równomiernie nakrytym},
a $V_\lambda$ \emph{arkuszami}.
Włókno $p^{-1}(x)$ jest dyskretne: arkusz zawierający
dany punkt włókna ma z tym włóknem tylko ten jeden
punkt wspólny.
\end{definicja}

\begin{przyklad}[Prosta nad okręgiem]
\label{prz:exp-nakrycie}
Odwzorowanie $p:\R\to S^1$, $p(t)=e^{2\pi i t}$,
jest nakryciem. Jeśli $U$ jest krótkim,
otwartym łukiem okręgu, wybierzmy przedział
$J$ długości mniejszej niż $1$, na którym
$p|_J$ jest homeomorfizmem na $U$.
Wtedy $p^{-1}(U)=\coprod_{k\in\Z}(J+k)$.
Włókno nad $1$ to $\Z$.
Podobnie $q_m:S^1\to S^1$, $q_m(z)=z^m$,
jest $m$-krotnym nakryciem dla $m\geq1$.
Rzeczywiście, jeśli
$U=\{e^{2\pi it}:a<t<b\}$ i $b-a<1$,
to jego przeciwobraz ma dokładnie $m$ arkuszy
$V_j=\{e^{2\pi i(t+j)/m}:a<t<b\}$,
$j=0,\ldots,m-1$. Na każdym z nich $q_m$
jest homeomorfizmem na $U$.
Dwukrotne nakrycie $\operatorname{Spin}(n)
\to\mathrm{SO}(n)$ pojawiło się już w
Twierdzeniu~\ref{thm:spin-nakrycie}.
\end{przyklad}

\begin{figure}[!htbp]
\centering
\begin{tikzpicture}[>=Stealth,line cap=round,line join=round,
 x={(1.55cm,0cm)},y={(.40cm,.37cm)},z={(0cm,1.25cm)}]
 % Okrąg u dołu, widoczna i zasłonięta połowa.
 \draw[manifoldblue!40,densely dashed,line width=1pt]
  plot[domain=0:180,samples=60,variable=\u]
   ({cos(\u)},{sin(\u)},0);
 \draw[manifoldblue,line width=1.8pt]
  plot[domain=180:360,samples=60,variable=\u]
   ({cos(\u)},{sin(\u)},0);
 % Włókno nad 1 oraz promień projekcji dla punktu t=1.25.
 \draw[manifoldblue!30,densely dotted]
  (1,0,0)--(1,0,2.9);
 \draw[manifoldgold!90!black,thick,dashed,
       -{Stealth[length=2.2mm]}]
  (0,1,1.65)--(0,1,0);
 \node[manifoldgold!70!black,right=4pt]
  at (0,1,.75) {$p$};
 % Fragment helisy: jeden obrót podnosi wysokość o jedną jednostkę.
 \draw[manifoldred,line width=2.1pt,
       -{Stealth[length=2.8mm]}]
  plot[domain=0:2.38,samples=180,variable=\t]
   ({cos(360*\t)},{sin(360*\t)},{.4+\t});
 \foreach \k in {0,1,2} {
  \fill[manifoldred] (1,0,{.4+\k}) circle(2.4pt);
  \node[manifoldred,right=5pt,font=\small]
   at (1,0,{.4+\k}) {$t=\k$};
 }
 \fill[manifoldblue] (1,0,0) circle(2.8pt);
 \node[manifoldblue,below right=3pt] at (1,0,0) {$1\in S^1$};
 \node[manifoldred!85!black,left=7pt]
  at (-1,0,2.7) {$\widetilde X\cong\R$};
 \node[manifoldblue!85!black,left=11pt]
  at (-1,0,-.22) {$X=S^1$};
\end{tikzpicture}
\caption{Nakrycie $p:\R\to S^1$, $p(t)=e^{2\pi i t}$,
przedstawione jako linia śrubowa nad okręgiem. Rzut pionowy
odpowiada $p$. Punkty linii śrubowej oznaczone $t=0,1,2$
leżą nad tym samym punktem $1\in S^1$.
Pokazano fragment linii śrubowej, która biegnie dalej
w obie strony.}
\label{fig:nakrycie-okragu}
\end{figure}

\begin{figure}[!htbp]
\centering
\begin{tikzpicture}[>=Stealth,line cap=round]
 \foreach \j in {0,1,2} {
  \pgfmathsetmacro{\yy}{1.35+0.75*\j}
  \pgfmathsetmacro{\xx}{-1.05+1.05*\j}
  \draw[manifoldgold!65!black,densely dashed,-{Stealth[length=2.2mm]}]
   (\xx,{\yy-.08}) -- (\xx,.44);
 }
 \foreach \j in {0,1,2} {
  \pgfmathsetmacro{\yy}{1.35+0.75*\j}
  \draw[manifoldblue,line width=2.1pt]
   (-1.7,\yy) .. controls (-.5,{\yy+.16}) and (.5,{\yy+.16}) .. (1.7,\yy);
  \node[manifoldblue,left=8pt] at (-1.7,\yy) {$V_{\j}$};
 }
 \draw[manifoldred,line width=2.3pt]
  (-1.7,.28) .. controls (-.5,.44) and (.5,.44) .. (1.7,.28);
 \node[manifoldred,left=8pt] at (-1.7,.28) {$U$};
 \node[manifoldgold!80!black,right=5pt] at (1.28,1.04) {$q_3$};
\end{tikzpicture}
\caption{Lokalny obraz nakrycia $q_3(z)=z^3$. Nad małym łukiem
$U\subset S^1$ znajdują się trzy rozłączne łuki $V_0,V_1,V_2$;
każdy odwzorowuje się homeomorficznie na $U$. Rysunek pokazuje
tylko tę lokalną sytuację: całkowita przestrzeń nakrycia $q_3$
jest jednym okręgiem.}
\label{fig:nakrycie-trzy-arkusze}
\end{figure}

W poniższych dowodach użyjemy \emph{liczby Lebesgue'a}:
dla otwartego pokrycia zwartej przestrzeni metrycznej istnieje
$\delta>0$ takie, że każdy niepusty zbiór o średnicy mniejszej
niż $\delta$ leży w jednym elemencie pokrycia.
Oto uzasadnienie. Dla każdego $x$ wybierzmy $\varepsilon_x>0$
tak, aby $B(x,2\varepsilon_x)$ mieściła się w elemencie pokrycia.
Ze zwartości skończenie wiele kul $B(x_i,\varepsilon_{x_i})$
pokrywa przestrzeń. Wystarczy $\delta=\min_i\varepsilon_{x_i}$:
jeśli $y$ należy do jednej z tych mniejszych kul, wszystkie
punkty zbioru zawierającego $y$ i o średnicy $<\delta$
leżą w odpowiadającej jej większej kuli. Stosujemy ten fakt
do odcinka i kwadratu z ich zwykłą metryką.

\begin{twierdzenie}[Podnoszenie drogi i homotopii]
\label{thm:podnoszenie-nakrycie}
Niech $p:E\to X$ będzie nakryciem.
\begin{enumerate}[label=(\alph*),leftmargin=*]
\item Każda droga $\gamma:I\to X$ i wybór
      $e_0\in p^{-1}(\gamma(0))$ wyznaczają
      dokładnie jedną drogę
      $\widetilde\gamma:I\to E$ z
      $\widetilde\gamma(0)=e_0$ i
      $p\widetilde\gamma=\gamma$.
\item Jeśli $H:I\times I\to X$ jest
      homotopią dróg, a $\widetilde H(s,0)$
      jest zadanym podniesieniem $H(s,0)$,
      to istnieje dokładnie jedno ciągłe
      podniesienie $\widetilde H$ całego $H$.
      Gdy $H$ zachowuje końce dróg,
      podniesione końce również są stałe.
\end{enumerate}
\end{twierdzenie}
\begin{proof}
\emph{Droga.} Przeciwobrazy równomiernie
nakrytych otwartych $U$ tworzą otwarte
pokrycie odcinka $I$. Z liczby Lebesgue'a
tego pokrycia wybieramy podział
$0=t_0<\cdots<t_N=1$ tak drobny, aby
każdy obraz $\gamma([t_{j-1},t_j])$
mieścił się w jednym $U_j$.
Arkusz nad $U_1$ zawierający $e_0$
określa na pierwszym przedziale wzór
$\widetilde\gamma(t)=(p|_{V_1})^{-1}(\gamma(t))$.
W punkcie $t_1$ wybieramy arkusz nad $U_2$
zawierający uzyskany punkt włókna i powtarzamy.
Wzory zgadzają się na końcach przedziałów,
więc sklejają się w drogę ciągłą.
Jeśli dwa podniesienia zgadzają się w pewnej
chwili, to lokalnie leżą w tym samym arkuszu
i są identyczne. Gdy w pewnej chwili
się różnią, leżą nad tym samym punktem
w różnych arkuszach, więc również w małej
okolicy chwili pozostają różne. Zbiór
chwil zgodności jest zatem otwarty
i domknięty w spójnym $I$ oraz
zawiera $0$, co daje jednoznaczność.

\emph{Homotopia.} Pokrycie $I^2$ zbiorami
$H^{-1}(U)$ ma liczbę Lebesgue'a.
Dzielimy kwadrat na skończoną siatkę
małych prostokątów, z których każdy ma
obraz w jednym równomiernie nakrytym $U$.
Na dolnej krawędzi mamy zadaną drogę.
Jej podniesienie na części dolnej krawędzi danego
prostokąta leży w jednym arkuszu nad wybranym $U$:
jest spójne, a arkusze są rozłączne i otwarte.
To samo dotyczy każdego wspólnego boku już podniesionego
prostokąta, gdy przechodzimy do następnego.
Podnosimy prostokąty kolejno od dołu i
od lewej: na każdym wybieramy arkusz
zawierający już ustalony obraz jego narożnika
i stosujemy $(p|_V)^{-1}\circ H$.
Na wspólnej krawędzi oba wzory są
podniesieniami tej samej drogi i zgadzają
się w narożniku, więc na mocy (a) są równe.
Skończone sklejanie daje ciągłą
$\widetilde H$; lokalna jednoznaczność
na spójnym $I^2$ daje jej jedyność.
Jeśli $H(0,t)$ albo $H(1,t)$ jest stała,
jej podniesienie leży w dyskretnym włóknie,
więc też jest stałe.
\end{proof}

Stąd $p_*:\pi_1(E,e_0)\to\pi_1(X,p(e_0))$
jest iniektywny. Jeśli obraz pętli jest
homotopijny do stałej pętli względem bazy,
podnosimy homotopię od danej pętli w $E$.
Końcowa podniesiona pętla leży w dyskretnym
włóknie i jest stała.

\begin{twierdzenie}[Kryterium podnoszenia odwzorowań]
\label{thm:kryterium-podnoszenia}
Niech $Y$ będzie łukowo spójna i lokalnie
łukowo spójna, $f:Y\to X$ ciągła,
$p:E\to X$ nakryciem, a $p(e_0)=f(y_0)$.
Istnieje ciągła $\widetilde f:Y\to E$
z $p\widetilde f=f$ i
$\widetilde f(y_0)=e_0$ wtedy i tylko wtedy, gdy
\[
 f_*\pi_1(Y,y_0)\subseteq
 p_*\pi_1(E,e_0).
\]
Takie podniesienie jest jedyne.
\end{twierdzenie}
\begin{proof}
Gdy $\widetilde f$ istnieje,
$f=p\widetilde f$, zatem
$f_*=p_*\widetilde f_*$, co daje konieczność.
Załóżmy inkluzję. Dla $y\in Y$ wybierzmy
drogę $\gamma$ od $y_0$ do $y$ i
zdefiniujmy $\widetilde f(y)$ jako koniec
podniesienia $f\circ\gamma$ z początkiem $e_0$.
Gdy $\gamma'$ jest innym wyborem,
pętla $f\circ(\gamma'*\bar\gamma)$
ma klasę w obrazie $p_*$. Jest więc
homotopijna względem punktu bazowego
do obrazu pewnej pętli w $E$ o początku
$e_0$. Podnosząc tę homotopię za pomocą
Twierdzenia~\ref{thm:podnoszenie-nakrycie},
otrzymujemy, że podniesienie
$f\circ(\gamma'*\bar\gamma)$ wraca do $e_0$.
Druga część tego zamkniętego podniesienia jest podniesieniem
$f\circ\bar\gamma$ kończącym się w $e_0$.
Po odwróceniu jest, z jednoznaczności, podniesieniem
$f\circ\gamma$ zaczynającym się w $e_0$.
Podniesienia $f\circ\gamma$ i $f\circ\gamma'$ mają więc ten sam koniec.

Sprawdźmy ciągłość, której nie daje sama
definicja punkt po punkcie. Wybierzmy
równomiernie nakryte $U$ zawierające $f(y)$
oraz łukowo spójne otwarte $W\ni y$ z
$f(W)\subseteq U$. Niech $V$ będzie
arkuszem nad $U$ zawierającym $\widetilde f(y)$.
Dla $z\in W$ połączmy $y$ z $z$ drogą
w $W$ i dopiszmy ją do drogi od $y_0$
do $y$. Definicja i podnoszenie drogi
pokazują wtedy, że
$\widetilde f(z)=(p|_V)^{-1}(f(z))$.
Na $W$ odwzorowanie jest więc ciągłe.
Jedyność wynika z jedyności podniesienia
każdej drogi od $y_0$.
\end{proof}

\begin{twierdzenie}[Grupa podstawowa okręgu]
\label{thm:pi1-okragu}
Podnoszenie względem $p:\R\to S^1$
z Przykładu~\ref{prz:exp-nakrycie} daje
izomorfizm
\[
 \operatorname{wind}:\pi_1(S^1,1)\xrightarrow{\ \cong\ }\Z,
 \qquad [\gamma]\longmapsto
 \widetilde\gamma(1),
 \quad \widetilde\gamma(0)=0.
\]
Liczba po prawej to \emph{liczba owinięć}.
\end{twierdzenie}
\begin{proof}
Ponieważ $\gamma(1)=1$, koniec
podniesienia leży w $p^{-1}(1)=\Z$.
Homotopijne pętle dają homotopię
podniesień z nieruchomymi końcami,
więc liczba zależy tylko od klasy.
Jeśli podniesienie $\alpha$ kończy się w $k$,
to podniesienie następnej pętli $\beta$
zaczynające się w $k$ jest przesunięciem
o $k$ podniesienia zaczynającego się w $0$.
Stąd $\operatorname{wind}([\alpha*\beta])
=\operatorname{wind}([\alpha])
+\operatorname{wind}([\beta])$.
Pętla $\gamma_k(s)=e^{2\pi i ks}$
ma podniesienie $ks$, więc homomorfizm
jest surjektywny. Gdy podniesienie
$\widetilde\gamma$ wraca do $0$,
wzór $(s,t)\mapsto(1-t)\widetilde\gamma(s)$
deformuje ją w $\R$ do stałej drogi
z zachowaniem końców. Projekcja tej
homotopii do $S^1$ deformuje $\gamma$
do pętli stałej. Jądro jest zerowe.
\end{proof}

Z Przykładu~\ref{prz:retrakcja-pierscien}
otrzymujemy $\pi_1(\R^2\setminus\{0\})\cong\Z$.
W szczególności okrąg i punkt nie są
homotopijnie równoważne. Obliczenie
$\pi_1(S^{n})=0$ dla $n\geq2$
wyniknie niżej z twierdzenia van Kampena.

\begin{przyklad}[Kiedy odwzorowanie okręgu podnosi się przez $q_m$?]
\label{prz:podnoszenie-poteg}
Niech $m\geq1$ i $f_k:S^1\to S^1$, $f_k(z)=z^k$,
gdzie $k\in\Z$. Z obliczenia liczby owinięć wynika, że
$(f_k)_*:\Z\to\Z$ jest mnożeniem przez $k$, a
$(q_m)_*$ jest mnożeniem przez $m$. Kryterium
podnoszenia mówi więc, że podniesienie $f_k$ przez
$q_m$, przyjmujące w punkcie $1$ wartość $1$, istnieje
dokładnie wtedy, gdy $k\Z\subseteq m\Z$, czyli $m\mid k$.
Gdy $k=m\ell$, podniesienie jest jawne:
$\widetilde f_k(z)=z^\ell$ i
$q_m\widetilde f_k(z)=z^{m\ell}=f_k(z)$.
Na przykład $z\mapsto z^6$ podnosi się przez $q_3$
jako $z\mapsto z^2$, natomiast $z\mapsto z^2$
nie ma ciągłego podniesienia przez $q_3$.
Przeszkodą jest liczba owinięć, a nie wybór lokalnej
gałęzi pierwiastka: taką gałąź można wybrać na każdym
krótkim łuku, lecz przy pełnym obiegu może ona wrócić
na inny arkusz.
\end{przyklad}

\begin{definicja}[Lokalna spójność i warunek półlokalny]
\index{lokalna spojnosc i warunek pollokalny@Lokalna spójność i warunek półlokalny}
\label{def:semilokalnie-prosto-spojna}
Przestrzeń $X$ jest \emph{lokalnie łukowo spójna},
gdy każdy punkt ma bazę otoczeń otwartych,
które są łukowo spójne. Jest
\emph{półlokalnie jednospójna}, gdy każdy
$x\in X$ ma otoczenie $U$, w którym
każda pętla oparta w $x$ jest homotopijna do stałej
względem punktu bazowego \emph{w całym $X$}. Nie wymagamy, aby
ta homotopia mieściła się w $U$.
Przestrzeń łukowo spójna z
$\pi_1(X,x_0)=0$ jest \emph{jednospójna}.
\end{definicja}

Każda rozmaitość topologiczna spełnia
lokalne warunki z tej definicji, gdyż
w otwartej mapie współrzędnych można
wybrać małą kulę. Warunek półlokalny
jest ważny: bez niego nakrycie uniwersalne
może w ogóle nie istnieć.

\begin{twierdzenie}[Konstrukcja nakrycia uniwersalnego]
\label{thm:uniwersalne-nakrycie}
Jeśli $X$ jest łukowo spójna, lokalnie
łukowo spójna i półlokalnie jednospójna,
to ma nakrycie $q:\widetilde X\to X$
z jednospójną przestrzenią $\widetilde X$.
\end{twierdzenie}
\begin{proof}
Ustalmy $x_0\in X$. Punktem $\widetilde X$
niech będzie klasa $[\gamma]$ drogi w $X$
zaczynającej się w $x_0$, przy czym
homotopie zachowują oba końce.
Definiujemy $q([\gamma])=\gamma(1)$.
Istnieje baza otwartych łukowo spójnych
zbiorów $U\subseteq X$, w których
klasa każdej pętli znika po inkluzji do $X$:
najpierw wybieramy otoczenie z warunku
półlokalnego, następnie w nim małe
łukowo spójne otoczenie z lokalnej
łukowej spójności. Każdą pętlę opartą w innym punkcie
takiego $U$ przenosimy do jego wyróżnionego punktu drogą w $U$;
izomorfizm zmiany punktu bazowego pokazuje, że ona też jest zerowa w $X$.

Dla drogi $\gamma$ kończącej się w $x\in U$
niech $W([\gamma],U)$ składa się z klas
$[\gamma*\eta]$, gdzie $\eta$ jest drogą
w $U$ od $x$ do dowolnego $y\in U$.
Jeśli dwie takie drogi $\eta,\eta'$
mają ten sam koniec, to
$\eta*\bar\eta'$ jest pętlą w $U$,
zerową w $X$. Zatem $[\gamma*\eta]
=[\gamma*\eta']$. Otrzymujemy bijekcję
$q:W([\gamma],U)\to U$. Bierzemy wszystkie
te zbiory jako bazę topologii. Sprawdźmy warunek przecięć:
jeśli $[\beta]\in W([\gamma],U)\cap W([\gamma'],U')$,
wybieramy otwarte łukowo spójne $O$ z naszej bazy,
zawierające $\beta(1)$ i zawarte w $U\cap U'$.
Wtedy $W([\beta],O)$ zawiera $[\beta]$ i mieści się w obu zbiorach:
do drogi kończącej się w $\beta(1)$ dopisujemy tylko drogę w $O$.
To dowodzi, że rzeczywiście otrzymujemy bazę.
Ograniczenie $q$ do każdego $W$ jest ciągłe, ponieważ przeciwobrazy
otwartych podzbiorów $U$ są sumami takich mniejszych zbiorów bazowych;
jest też otwarte, bo mniejsze $W([\beta],O)$ przechodzą na otwarte $O$.
Wcześniej wykazana bijekcja jest więc homeomorfizmem.
Jeśli dwa $W$ nad tym samym
$U$ mają wspólny punkt, ich drogi początkowe
różnią się drogą w $U$ z dokładnością
do homotopii w $X$, więc zbiory $W$
są równe. Dla mniejszych otoczeń
używamy ograniczenia tych samych odwzorowań.
Stąd $q^{-1}(U)$ jest rozłączną sumą
arkuszy $W([\gamma],U)$: $q$ jest nakryciem.

Jest ono łukowo spójne. Drogę od klasy
stałej $[c_{x_0}]$ do $[\gamma]$
otrzymujemy, biorąc po kolei klasy
odcinków $\gamma|_{[0,t]}$ po
przeparametryzowaniu na $I$; lokalnie
leżą one w jednym $W$, więc tworzą
ciągłą drogę. Jeśli pętla
$\lambda$ w $\widetilde X$ jest oparta
w $[c_{x_0}]$, jej rzut $\gamma=q\lambda$
ma podniesienie $\lambda$ i kończy
się w tej samej klasie drogi.
Konstrukcja podniesienia pokazuje,
że ta klasa jest $[\gamma]$, więc
$\gamma$ jest homotopijna do stałej
pętli względem bazowego końca.
Podnosimy tę homotopię, zaczynając
od $\lambda$, za pomocą
Twierdzenia~\ref{thm:podnoszenie-nakrycie}.
Na końcu dostajemy stałą pętlę
w dyskretnym włóknie. Zatem
$\pi_1(\widetilde X)=0$.
\end{proof}

\begin{uwaga}[Skąd podgrupy grupy podstawowej?]
\label{uwaga:nakrycia-podgrupy}
Dla spójnego nakrycia $p:(E,e_0)\to(X,x_0)$
iniekcja $p_*$ wyróżnia podgrupę
$H=p_*\pi_1(E,e_0)\leq\pi_1(X,x_0)$.
W warunkach Twierdzenia~\ref{thm:uniwersalne-nakrycie}
każda taka podgrupa daje nakrycie:
na klasach dróg w $\widetilde X$ grupa
$\pi_1(X,x_0)$ działa przez
$[\omega]\cdot[\gamma]=[\omega*\gamma]$,
a iloraz $H\backslash\widetilde X$
jest nakryciem $X$. Istotnie, działanie przenosi
$W([\gamma],U)$ na $W([\omega*\gamma],U)$,
więc jest działaniem przez homeomorfizmy. Jest wolne:
z $[\omega*\gamma]=[\gamma]$ po dopisaniu $\bar\gamma$
wynika $[\omega]=1$.
Co więcej, jeśli $hW\cap W\ne\varnothing$, to $hW=W$,
bo są to arkusze nad tym samym $U$. Ponieważ $q$ jest
iniektywne na $W$ i $q(hz)=q(z)$, element $h$ ustala punkt $z\in W$;
stąd $h=1$. Różne przesunięcia jednego arkusza są więc rozłączne.
Projekcja na orbity jest otwarta: przeciwobraz obrazu zbioru
otwartego $O$ jest sumą otwartych zbiorów $hO$.
Obrazy arkuszy w ilorazie są zatem otwarte i każdy odwzorowuje się
homeomorficznie na $U$. To uzasadnia własność nakrycia.
Przestrzeń ilorazowa jest łukowo spójna jako obraz $\widetilde X$.
Podniesienie pętli $\gamma$ z $x_0$ kończy się w orbicie $[\gamma]$;
wraca do orbity drogi stałej dokładnie wtedy, gdy $[\gamma]\in H$.
Ponieważ obraz homomorfizmu indukowanego przez nakrycie składa się
właśnie z klas mających zamknięte podniesienie, skonstruowane nakrycie
realizuje dokładnie podgrupę $H$.
Z drugiej strony kryterium
\ref{thm:kryterium-podnoszenia} podnosi
$q:\widetilde X\to X$ do $E$, z ustaloną wartością $e_0$.
Przy obecnych założeniach $E$ jest lokalnie łukowo spójna,
więc jej spójność daje łukową spójność: składowe łukowe są otwarte.
Każdy punkt $E$ można połączyć z $e_0$ drogą; rzut tej drogi
pokazuje, że podniesione odwzorowanie $\widetilde X\to E$ jest surjektywne.
Dwa punkty $[\gamma],[\gamma']$ nad tym samym punktem $X$
trafiają w ten sam punkt $E$
dokładnie wtedy, gdy $[\gamma'*\bar\gamma]\in H$.
W jedną stronę składamy podniesione drogi o wspólnym końcu,
otrzymując pętlę w $E$; w drugą używamy podnoszenia homotopii
jak w dowodzie kryterium podnoszenia. Ten warunek jest też
równoważny przynależności do jednej orbity, gdyż
$[\gamma'*\bar\gamma]\cdot[\gamma]=[\gamma']$.
Wybierzmy spójne, równomiernie nakryte otoczenie $U$
i arkusz $W$ nakrycia $q$ nad $U$. Obraz $W$ przez
podniesione odwzorowanie jest spójny, więc mieści się
w jednym arkuszu $V$ nakrycia $p$ nad $U$.
Na $W$ odwzorowanie to ma postać
$(p|_V)^{-1}\circ q|_W$, a zatem jest homeomorfizmem
$W\to V$. Wykazana wyżej zgodność włókien z orbitami
$H$ pokazuje, że po przejściu do ilorazu otrzymujemy
lokalnie homeomorfizm. To daje
$E\cong H\backslash\widetilde X$
nad $X$. Zmiana wybranego punktu
$e_0$ sprzęga podgrupę $H$: droga w $E$ między dwoma punktami
nad $x_0$ rzutuje się na pętlę $u$, a zmiana bazy zastępuje
$H$ podgrupą $[\bar u]H[u]$.
Odwrotnie, dowolną pętlę $u$ w $X$ można podnieść z $e_0$;
jej koniec wybiera punkt włókna realizujący właśnie to sprzężenie.
Tak powstaje klasyfikacja spójnych
nakryć przez klasy sprzężoności
podgrup $\pi_1(X,x_0)$.
Ponadto włókno nad $x_0$ jest w tej konstrukcji
zbiorem lewych warstw $H\backslash\pi_1(X,x_0)$.
Liczba arkuszy spójnego nakrycia wynosi więc
$[\pi_1(X,x_0):H]$, także gdy indeks jest nieskończony.
Droga między dwoma punktami bazy daje przez podnoszenie
bijekcję ich włókien; droga odwrotna daje bijekcję odwrotną.
\end{uwaga}

\begin{definicja}[Transformacje nakrycia i nakrycie normalne]
\index{transformacja nakrycia}
\label{def:transformacja-nakrycia}
\emph{Transformacją nakrycia} $p:E\to X$ nazywamy
homeomorfizm $T:E\to E$ spełniający $pT=p$.
Takie homeomorfizmy tworzą grupę $\operatorname{Deck}(p)$.
Nakrycie spójne nazywamy \emph{normalnym} (lub regularnym),
jeśli transformacje nakrycia działają tranzytywnie na każdym
włóknie: dowolne dwa punkty leżące nad tym samym punktem
bazy można zamienić pewną transformacją. Grupa ta opisuje
symetrie przestrzeni całkowitej, których nie widać w bazie.
\end{definicja}

\begin{twierdzenie}[Grupa transformacji nakrycia]
\label{thm:deck-normalizator}
Niech $X$ będzie łukowo spójna, lokalnie łukowo spójna
i półlokalnie jednospójna, a $p:(E,e_0)\to(X,x_0)$
spójnym nakryciem. Dla
$G=\pi_1(X,x_0)$ i $H=p_*\pi_1(E,e_0)$ zachodzi
\[
 \operatorname{Deck}(p)\cong N_G(H)/H,
 \qquad N_G(H)=\{g\in G:gHg^{-1}=H\}.
\]
Nakrycie jest normalne wtedy i tylko wtedy, gdy
$H\mathrel{\triangleleft}G$; wówczas
$\operatorname{Deck}(p)\cong G/H$.
\end{twierdzenie}
\begin{proof}
Z klasyfikacji powyżej możemy utożsamić $E$
z $H\backslash\widetilde X$. Dla $g\in N_G(H)$
wybierzmy reprezentującą go pętlę i zdefiniujmy
$T_g(H[\gamma])=H[g*\gamma]$; wartość zależy tylko
od klasy $g$, ponieważ homotopie pętli zachowują
punkt bazowy.
Jeżeli $H[\gamma]=H[h*\gamma]$ dla $h\in H$,
to w grupie $G$ mamy $gh=(ghg^{-1})g$,
a $ghg^{-1}\in H$.
Zatem wzór nie zależy od reprezentanta orbity.
Jest ciągły, ponieważ pochodzi z działania $g$
na $\widetilde X$; odwrotnością jest $T_{g^{-1}}$.
Ponieważ pętla $g$ nie zmienia końca drogi $\gamma$,
otrzymujemy $pT_g=p$. Równość
$T_gT_{g'}=T_{gg'}$ pokazuje, że jest to
homomorfizm $N_G(H)\to\operatorname{Deck}(p)$.
Jego jądrem jest $H$: transformacja $T_g$
ustala punkt $H[c_{x_0}]$ dokładnie wtedy, gdy $g\in H$,
a taka transformacja, jako podniesienie $p$,
jest wyznaczona przez wartość w jednym punkcie.

Pokażmy surjektywność. Dla dowolnej transformacji $T$
połączmy $e_0$ z $T(e_0)$ drogą $\alpha$ w $E$.
Jej rzut jest pętlą $u$ w $X$. Zmiana punktu bazowego
wzdłuż $\alpha$ daje
$p_*\pi_1(E,T(e_0))=[\bar u]H[u]$.
Z drugiej strony $pT=p$, więc ta grupa jest równa $H$.
Stąd $[u]\in N_G(H)$. Transformacje $T_u$ i $T$
przenoszą $e_0$ w ten sam punkt włókna; jedyność
podnoszenia dróg i łukowa spójność $E$ dają $T_u=T$.
Twierdzenie o izomorfizmie grup kończy pierwszy dowód.

Włókno nad $x_0$ odpowiada warstwom $H\backslash G$,
a orbita punktu $H$ pod działaniem transformacji
odpowiada warstwom reprezentowanym przez $N_G(H)$.
Tranzytywność jest więc równoważna $N_G(H)=G$, czyli
normalności $H$. Jeśli zachodzi dla jednego włókna,
zachodzi dla każdego: podnoszenie drogi z $x_0$ do
innego punktu bazy przenosi oba włókna bijektywnie
i komutuje z każdą transformacją nakrycia.
\end{proof}

\begin{przyklad}[Symetrie nakryć okręgu]
Dla $p:\R\to S^1$ transformacje to przesunięcia
$T_n(t)=t+n$, $n\in\Z$. Istotnie, każda transformacja
jest wyznaczona przez obraz $0$, który musi leżeć
w $p^{-1}(1)=\Z$, a przesunięcia realizują wszystkie
te możliwości. Zatem $\operatorname{Deck}(p)\cong\Z$.
Dla $q_m(z)=z^m$ transformacje są obrotami
$T_j(z)=e^{2\pi i j/m}z$, $j=0,\ldots,m-1$,
więc $\operatorname{Deck}(q_m)\cong\Z/m\Z$.
To zgadza się z podgrupą $m\Z\leq\pi_1(S^1)=\Z$:
ma ona indeks $m$ i jest normalna. Ponieważ wszystkie
podgrupy $\Z$ mają postać $m\Z$ dla $m\geq0$,
każde spójne nakrycie okręgu jest, z dokładnością
do izomorfizmu nad $S^1$, jednym z $q_m$ dla $m\geq1$
albo nakryciem $p$ odpowiadającym podgrupie $0\Z=\{0\}$.
\end{przyklad}

\section{Twierdzenie van Kampena i pierwsze obliczenia}
\label{sec:van-kampen}

Nakrycie okręgu mierzy jedną liczbę owinięć.
Gdy przestrzeń składa się z kilku części,
potrzebujemy sposobu łączenia informacji
o pętlach w tych częściach.
Ważne są przy tym dwie rzeczy: pętle porównujemy
przy tym samym punkcie bazowym, a pętla leżąca
w obu częściach nie może dawać dwóch niezależnych
generatorów. Twierdzenie van Kampena ujmuje
drugą zasadę jako relację w iloczynie wolnym.

\begin{twierdzenie}[van Kampen: dwa otwarte zbiory]
\label{thm:van-kampen}
Niech $X=U\cup V$, gdzie $U,V$ są otwarte,
łukowo spójne, a $U\cap V$ jest łukowo
spójne i zawiera $x_0$. Jeśli
$G_U=\pi_1(U,x_0)$, $G_V=\pi_1(V,x_0)$
oraz $G_0=\pi_1(U\cap V,x_0)$, a
$i_U:U\cap V\hookrightarrow U$ i $i_V:U\cap V\hookrightarrow V$
są inkluzjami, to
\[
 \pi_1(X,x_0)\cong
 (G_U*G_V)/
 \left\langle\!\left\langle
 (i_U)_*(h)(i_V)_*(h)^{-1}:h\in G_0
 \right\rangle\!\right\rangle.
\]
Iloczyn wolny wprowadziliśmy w Definicji~\ref{def:free-groups-5};
postać normalną i własność uniwersalną uzasadnia
Stwierdzenie~\ref{prop:free-product-normal-5}. Zapis
$\langle\!\langle\cdot\rangle\!\rangle$
oznacza domknięcie normalne z
Definicji~\ref{def:normal-quotient-5}.
W szczególności gdy $U\cap V$ jest po prostu spójne,
nie ma relacji pochodzących z przecięcia i otrzymujemy
$\pi_1(X,x_0)\cong G_U*G_V$.
\end{twierdzenie}
\begin{proof}
Inkluzje $U,V\hookrightarrow X$, przez własność uniwersalną
iloczynu wolnego, określają homomorfizm
$\Phi:G_U*G_V\to\pi_1(X,x_0)$.
Pętla w $U\cap V$ daje ten sam obraz
niezależnie od wyboru strony.
Zatem wszystkie relacje z tezy
leżą w $\ker\Phi$.

\emph{Surjektywność.} Dla pętli
$\alpha:I\to X$ przeciwobrazy $U,V$
tworzą otwarte pokrycie zwartego odcinka.
Z lematu o liczbie Lebesgue'a dla tego pokrycia
dzielimy odcinek na tyle małych przedziałów,
aby obraz każdego leżał w $U$ albo $V$.
Sąsiednie kawałki przypisane temu
samemu zbiorowi łączymy. Każdy
punkt przejścia między kawałkiem
w $U$ i kawałkiem w $V$ leży w
$U\cap V$. Łączymy go z $x_0$
pomocniczą drogą w $U\cap V$.
Jeśli $\lambda_j$ prowadzi od $x_0$
do $j$-tego punktu przejścia,
to z kawałka $\alpha_j$ budujemy
pętlę $\lambda_{j-1}*\alpha_j*
\bar\lambda_j$ całkowicie w $U$
albo w $V$. Iloczyn tych pętli
odtwarza $\alpha$: sąsiednie
$\bar\lambda_j*\lambda_j$
kasują się na mocy
Twierdzenia~\ref{thm:grupa-podstawowa}.
Na początku i końcu bierzemy stałą drogę
$\lambda_0=\lambda_m=c_{x_0}$, więc nie pojawia się
dodatkowe sprzężenie całej pętli.

Zmiana drogi $\lambda_j$ na $\lambda'_j$
wyznacza pętlę $h=\lambda'_j*\bar\lambda_j$ w $U\cap V$
\emph{opartą w $x_0$}. Poprzedni czynnik słowa mnożymy z prawej
przez $h^{-1}$, a następny z lewej przez $h$.
Oba zapisy $h$ są równe w ilorazie z tezy, więc się skracają.
Również dopisanie podziału wewnątrz
jednego otwartego zbioru nie zmienia
otrzymanej klasy słowa.

\emph{Jądro.} Przypuśćmy, że słowo
z $G_U*G_V$ daje pętlę homotopijną
do stałej w $X$. Wypełnijmy ją
homotopią $F:I^2\to X$, stałą na
pozostałych trzech bokach.
Zwartość kwadratu pozwala wybrać
siatkę tak drobną, by obraz
każdego małego prostokąta leżał
w całości w $U$ lub w $V$.
Podział dolnego boku dobieramy także tak, aby zawierał końce
pętli reprezentujących litery wyjściowego słowa.
Dla każdego wierzchołka $z$ siatki wybierzmy drogę $\mu_z$
od $x_0$ do $F(z)$. Jeśli $F(z)\in U\cap V$, wybieramy ją
w przecięciu; jeśli $F(z)$ należy tylko do jednego zbioru,
wybieramy drogę w tym zbiorze. Jest to możliwe dzięki założeniom
o łukowej spójności. Przy $F(z)=x_0$ wybieramy drogę stałą.

Zorientowanej krawędzi $e:z\to w$ przypisujemy klasę
$[\mu_z*(F|_e)*\bar\mu_w]$ w grupie zbioru zawierającego
obraz sąsiedniego prostokąta. Obie drogi pomocnicze leżą w tym
zbiorze. Jeżeli dwa sąsiednie prostokąty mają różne oznaczenia,
ich wspólna krawędź i obie drogi pomocnicze leżą w $U\cap V$,
więc przypisania dają ten sam element ilorazu z tezy.
Odwrócenie krawędzi daje element odwrotny.

Iloczyn czterech elementów wokół jednego prostokąta jest jedynką:
drogi pomocnicze skracają się parami, a jego brzeg jest homotopijny
do stałej we właściwym $U$ lub $V$, przez odwzorowanie $F$ na tym prostokącie.
Możemy zatem zastępować dwa kolejne boki prostokąta pozostałymi
dwoma, z tymi samymi końcami, bez zmiany elementu ilorazu.
Przesuwając tak łamaną po kolejnych prostokątach, wiersz po wierszu,
zamieniamy dolny bok kwadratu na drogę wzdłuż pozostałych trzech boków.
Te trzy boki są stałe. Na dolnym boku drogi pomocnicze wewnątrz
każdej pierwotnej litery skracają się w jej grupie.
Rzeczywiście, jeżeli litera leży w $U$, ale prostokąt
przylegający do fragmentu dolnego boku ma etykietę $V$,
ten fragment i jego końce leżą w $U\cap V$; relacja
z przecięcia pozwala więc przypisać mu etykietę $U$.
Analogicznie postępujemy dla litery w $V$.
Przy końcach liter drogi $\mu_z$ są stałe, ponieważ
obraz tych wierzchołków jest równy $x_0$.
Odtwarzamy więc dokładnie wyjściowe słowo,
które okazało się jedynką w ilorazie. Każdy krok używał tylko
relacji wewnątrz obu grup i relacji pochodzących z przecięcia.
Stąd $\ker\Phi$ zawiera się we wskazanym domknięciu normalnym
relacji $N$; przeciwna inkluzja wynika z pierwszego kroku,
bo jądro jest podgrupą normalną. Wzór $gN\mapsto\Phi(g)$
jest więc dobrze określonym homomorfizmem na ilorazie.
Jest surjektywny, a jego jądro jest trywialne, ponieważ
$\Phi(g)=1$ zachodzi dokładnie dla $g\in N$. Daje to żądany izomorfizm.
\end{proof}

\begin{przyklad}[Dwa okręgi złączone w punkcie]
\label{prz:pi1-dwa-okregi}
Niech $X=S^1_a\vee S^1_b$.
Wybierzmy na każdym okręgu po jednym
punkcie różnym od punktu sklejenia
$x_0$. Usuwając punkt z drugiego
okręgu, otrzymujemy otwarty $U$
retrakcyjny na $S^1_a$; analogicznie
$V$ retrakcyjny na $S^1_b$.
Przecięcie $U\cap V$ jest ściągalne
(to dwa otwarte łuki spotykające się
w $x_0$). Wobec tego
$\pi_1(S^1\vee S^1,x_0)\cong\Z*\Z=:F(a,b)$.
Słowa $ab$ i $ba$ są różne:
grupa podstawowa nie musi być przemienna.
\end{przyklad}

\begin{przyklad}[Wypełnienie dziury relacją]
\label{prz:van-kampen-dysk}
Niech $X=D^2=\{y\in\R^2:\|y\|\leq1\}$,
$U=D^2\setminus\{0\}$, $V=\{y:\|y\|<1/2\}$
oraz $x_0=(1/4,0)$. Są to otwarte podzbiory $X$,
a $U\cap V$ jest przekłutym małym dyskiem.
Zarówno $U$, jak i $U\cap V$ deformują się radialnie
na okrąg o promieniu $1/4$: dla $y\ne0$ wystarcza
homotopia $H(y,t)=((1-t)\|y\|+t/4)y/\|y\|$.
Promień pozostaje między $\|y\|$ a $1/4$,
więc homotopia nie opuszcza żadnego z tych zbiorów.
Ich grupy podstawowe są zatem równe $\Z$.
Inkluzja $U\cap V\hookrightarrow U$
indukuje izomorfizm, bo obie retrakcje ustalają
ten sam okrąg. Zbiór $V$ jest wypukły, więc
$\pi_1(V,x_0)=0$. W twierdzeniu van Kampena
generator $a$ pochodzący z $U$ jest zatem równy
jedynce po stronie $V$. Otrzymujemy
$\pi_1(D^2,x_0)\cong\Z/\langle a\rangle=0$.
Geometrycznie dysk $V$ wypełnia pętlę obiegającą
brakujący w $U$ punkt.
\end{przyklad}

\begin{figure}[!htbp]
\centering
\begin{tikzpicture}[>=Stealth,line cap=round,line join=round]
 \fill[manifoldblue!12] (0,0) circle(2cm);
 \fill[manifoldgreen!18] (0,0) circle(1cm);
 \draw[manifoldblue,thick] (0,0) circle(2cm);
 \draw[manifoldgreen!65!black,thick,dashed] (0,0) circle(1cm);
 \draw[manifoldred,thick,-{Stealth[length=2.2mm]}]
  (.5,0) arc[start angle=0,end angle=330,radius=.5cm];
 \fill[manifoldgreen!65!black] (0,0) circle(3pt);
 \fill[manifoldred] (.5,0) circle(2.5pt);
 \node[manifoldgreen!65!black,below left=2pt] at (0,0) {$0$};
 \node[right=4pt] at (.5,0) {$x_0$};
 \node[manifoldblue!80!black] at (1.48,.7) {$U$};
 \node[manifoldgreen!55!black] at (-.53,.55) {$V$};
 \node[manifoldred] at (-.7,-.36) {$a$};
\end{tikzpicture}
\caption{Pokrycie dysku przez przekłuty dysk $U$ i mały dysk $V$.
Zielony punkt $0$ należy do $V$, lecz nie do $U$.
Pętla $a$ w przecięciu generuje $\pi_1(U)$,
lecz kurczy się w $V$; relacja van Kampena ją usuwa.}
\label{fig:van-kampen-dysk}
\end{figure}

\begin{uwaga}[Dlaczego przecięcie musi być łukowo spójne]
Pokryjmy $S^1\subset\C$ zbiorami
$U=S^1\setminus\{-1\}$ i $V=S^1\setminus\{1\}$,
wybierając $x_0=i$. Każdy z nich jest otwartym łukiem,
więc ma trywialną grupę podstawową. Natomiast
$U\cap V$ ma dwie składowe łukowe.
Gdyby bez sprawdzenia hipotezy zastosować wzór
z Twierdzenia~\ref{thm:van-kampen}, otrzymalibyśmy
grupę trywialną, chociaż $\pi_1(S^1,x_0)\cong\Z$
na mocy Twierdzenia~\ref{thm:pi1-okragu}.
Druga składowa przecięcia pozwala przejść z jednego
łuku do drugiego i wrócić, tworząc generator okręgu.
\end{uwaga}

\begin{przyklad}[Sfera wymiaru co najmniej dwa]
\label{prz:pi1-sfera}
Niech $n\geq2$. Zbiory
$U=S^n\setminus\{N\}$ i
$V=S^n\setminus\{S\}$, gdzie $N,S$
są biegunami, są homeomorficzne
z $\R^n$, więc mają trywialne
grupy podstawowe. Ich przecięcie
jest łukowo spójne: deformuje się
na równik $S^{n-1}$, który jest
łukowo spójny dla $n-1\geq1$.
Twierdzenie~\ref{thm:van-kampen}
daje $\pi_1(S^n)=0$.
Stąd $\pi_1(\R^{n+1}\setminus\{0\})=0$
dla $n\geq2$, co odróżnia ten
przypadek od płaszczyzny.
\end{przyklad}

\section{Sympleksy i triangulacje}
\label{sec:sympleksy}

Sympleks jest najprostszym wielowymiarowym
klockiem: odcinek w wymiarze $1$, trójkąt
w wymiarze $2$, czworościan w wymiarze $3$.
Współrzędne barycentryczne opisują
zarazem punkt i to, na której ścianie leży.

\begin{definicja}[Sympleks standardowy i ściany]
\index{sympleks standardowy i sciany@Sympleks standardowy i ściany}
\label{def:sympleks-standardowy}
Dla $n\geq0$ \emph{standardowy $n$-sympleks}
to
\[
 \Delta^n=
 \left\{(t_0,\ldots,t_n)\in\R^{n+1}:
 t_i\geq0,\ \sum_{i=0}^{n}t_i=1\right\}.
\]
Jego wierzchołki to wektory $e_0,\ldots,e_n$.
\emph{Wnętrze} oznacza $t_i>0$ dla wszystkich $i$
(względem afinicznej otoczki sympleksu).
Dla $n\geq1$ \emph{ściana kowymiaru jeden} naprzeciw $e_i$
jest zadana równaniem $t_i=0$ i jest kopią $\Delta^{n-1}$.
Ogólniej ścianę wyznacza dowolny niepusty podzbiór wierzchołków;
jest właściwa, gdy pomija choć jeden z nich.
Suma wszystkich właściwych ścian to $\partial\Delta^n$.
Przy $n=0$ nie ma niepustych właściwych ścian i $\partial\Delta^0=\varnothing$.
Punkty $v_0,\ldots,v_n\in\R^N$
są \emph{afinicznie niezależne}, jeśli
$v_1-v_0,\ldots,v_n-v_0$ są liniowo
niezależne. Ich otoczka wypukła
$[v_0,\ldots,v_n]$ jest
\emph{geometrycznym $n$-sympleksem}.
\end{definicja}

Każdy punkt sympleksu ma postać
$x=\sum_i t_iv_i$. Współczynniki
$t_i$ są jednoznaczne: jeśli
$\sum_i(t_i-s_i)v_i=0$ i obie sumy
współczynników są równe $1$, to
$\sum_{i=1}^n(t_i-s_i)(v_i-v_0)=0$;
afiniczna niezależność wymusza
$t_i=s_i$ dla wszystkich $i$.
W szczególności usunięcie jednego
wierzchołka daje ścianę
o wymiarze $n-1$.
Odwzorowanie $\Delta^n\to[v_0,\ldots,v_n]$ dane wzorem
$(t_0,\ldots,t_n)\mapsto\sum_i t_iv_i$ jest homeomorfizmem:
jest ciągłą bijekcją, a jej dziedzina jest zwarta,
zaś obraz jest podzbiorem przestrzeni Hausdorffa.
Współrzędne barycentryczne pozwalają więc przenosić
rachunki ze standardowego sympleksu na dowolny
sympleks geometryczny.

\begin{figure}[!htbp]
\centering
\begin{tikzpicture}[line cap=round,line join=round]
 \node[font=\bfseries] at (0,3.45) {$\Delta^0$};
 \node[font=\bfseries] at (4.05,3.45) {$\Delta^1$};
 \node[font=\bfseries] at (8.63,3.45) {$\Delta^2$};
 % Wymiar zero.
 \fill[manifoldblue!10] (0,1.15) circle(.26);
 \fill[manifoldred] (0,1.15) circle(3.2pt);
 \node[below=9pt] at (0,1.15) {$v_0$};
 % Wymiar jeden.
 \coordinate (b0) at (2.75,1.15);
 \coordinate (b1) at (5.35,1.15);
 \draw[manifoldblue,line width=2.6pt] (b0)--(b1);
 \fill[manifoldred] (b0) circle(3.2pt)
  (b1) circle(3.2pt);
 \node[below=8pt] at (b0) {$v_0$};
 \node[below=8pt] at (b1) {$v_1$};
 % Wymiar dwa.
 \coordinate (c0) at (7.00,.35);
 \coordinate (c1) at (10.28,.35);
 \coordinate (c2) at (8.21,2.45);
 \fill[manifoldblue!13] (c0)--(c1)--(c2)--cycle;
 \draw[manifoldblue,line width=2pt] (c0)--(c1)--(c2)--cycle;
 \fill[manifoldred] (c0) circle(3.2pt)
  (c1) circle(3.2pt) (c2) circle(3.2pt);
 \node[below left=3pt] at (c0) {$v_0$};
 \node[below right=3pt] at (c1) {$v_1$};
 \node[above=4pt] at (c2) {$v_2$};
\end{tikzpicture}
\caption{Sympleksy wymiarów $0,1,2$: punkt,
odcinek i wypełniony trójkąt. Wymiar
$n$ odpowiada $n+1$ afinicznie niezależnym
wierzchołkom; ściany powstają przez ich usunięcie.}
\label{fig:sympleks-trojkat}
\end{figure}

\begin{uwaga}[Brzeg geometryczny i znaki orientacji]
\label{uwaga:orientacja-przed-homologia}
Jako zbiór sympleks $[v_0,v_1,v_2]$ nie zależy
od kolejności wypisania wierzchołków.
Gdy ten sam zapis oznacza \emph{zorientowany łańcuch},
zamiana dwóch wierzchołków zmienia znak.
Na przykład algebraiczny brzeg zorientowanego trójkąta to
\[
 \partial[v_0,v_1,v_2]
 =[v_1,v_2]-[v_0,v_2]+[v_0,v_1].
\]
Środkowa krawędź jest obiegana od $v_2$ do $v_0$,
a jej wybrany kierunek od $v_0$ do $v_2$ daje znak minus.
Natomiast $\partial\Delta^2$ w Definicji~\ref{def:sympleks-standardowy}
jest po prostu \emph{zbiorem} trzech krawędzi, bez znaków.
Definicja łańcuchów i dowód reguły $\partial^2=0$
pojawią się w~\S\ref{sec:homologia-symplicjalna}.
\end{uwaga}

\begin{stwierdzenie}[Brzeg sympleksu jest sferą]
\label{prop:brzeg-sympleksu}
Dla $n\geq1$ przestrzeń $\Delta^n$
jest homeomorficzna z dyskiem $D^n$,
a $\partial\Delta^n$ z $S^{n-1}$.
\end{stwierdzenie}
\begin{proof}
Pracujemy w $n$-wymiarowej płaszczyźnie
afinicznej $P=\{\sum t_i=1\}$.
Środek $c=(1/(n+1),\ldots,1/(n+1))$
leży we wnętrzu. Każdy promień z $c$
w $P$ przecina brzeg dokładnie raz:
nierówności $t_i\geq0$ wyznaczają
na nim odcinek od $c$ do pierwszej
ściany, na której pewna współrzędna
staje się zerem.
Przypisanie punktowi brzegowemu
$x$ kierunku $(x-c)/\|x-c\|$
jest ciągłą bijekcją
$\partial\Delta^n\to S^{n-1}$
po utożsamieniu $P-c\cong\R^n$.
Dziedzina jest zwarta, a sfera
Hausdorffa, więc odwrotność jest
ciągła. Jeśli $b(u)$ jest tą odwrotnością dla kierunku
$u\in S^{n-1}$, wzór $ru\mapsto c+r(b(u)-c)$, $0<r\leq1$,
z wartością $c$ w zerze daje ciągłą bijekcję z $D^n$ na $\Delta^n$.
Ciągłość w zerze wynika z ograniczoności $b(u)-c$.
Ponownie zwartość i własność Hausdorffa dają homeomorfizm,
który na brzegu jest poprzednią identyfikacją. Dla $n=1$ brzeg
składa się z dwóch punktów,
czyli jest $S^0$.
\end{proof}

\begin{definicja}[Skończony kompleks symplicjalny]
\index{skonczony kompleks symplicjalny@Skończony kompleks symplicjalny}
\label{def:kompleks-symplicjalny}
Niech $V$ będzie skończonym zbiorem
wierzchołków. \emph{Kompleks symplicjalny}
$K$ to rodzina niepustych skończonych
podzbiorów $V$ taka, że jeśli
$\sigma\in K$ i $\varnothing\ne\tau\subseteq\sigma$,
to $\tau\in K$. Zbiór o $n+1$
wierzchołkach nazywamy
\emph{$n$-sympleksem} $K$.
Jego \emph{realizacja geometryczna} to
\[
 |K|=\left\{\sum_{v\in V}t_v e_v:
 t_v\geq0,\ \sum_v t_v=1,\
 \{v:t_v>0\}\in K\right\}\subseteq\R^{V}.
\]
Kompleks $L\subseteq K$ zamknięty
na przechodzenie do ścian jest
\emph{podkompleksem}.
\emph{Skończona triangulacja} przestrzeni $X$
to homeomorfizm $|K|\to X$. W tej sekcji ograniczamy się do
skończonych kompleksów; ich realizacje są zwarte. Triangulacje
przestrzeni niezwartych wymagają dopuszczenia nieskończenie wielu sympleksów.
\end{definicja}

Zwartość $|K|$ wynika też bezpośrednio z definicji:
jest ono skończoną unią domkniętych sympleksów
leżących w standardowym sympleksie na zbiorze $V$.

W tej definicji dwie różne ściany
mają dokładnie to wspólne, co
wynika ze wspólnych wierzchołków.
Istotnie, punkt części wspólnej sympleksów
na zbiorach wierzchołków $\sigma$ i $\tau$ ma,
dzięki jednoznaczności współrzędnych barycentrycznych,
dodatnie współczynniki tylko przy wierzchołkach
z $\sigma\cap\tau$. Część wspólna jest zatem
ich wspólną ścianą albo jest pusta.
Słowo \emph{podkompleks} jest
ważniejsze niż dowolny podzbiór:
zawiera wraz z sympleksem
wszystkie jego ściany.

\begin{przyklad}[Okrąg i dysk z tych samych wierzchołków]
Niech $V=\{v_0,v_1,v_2\}$, a $K_{\partial}$
zawiera wszystkie niepuste podzbiory $V$
o jednym lub dwóch elementach. Jego realizacja
składa się z trzech wierzchołków i trzech krawędzi,
więc $|K_{\partial}|=\partial\Delta^2\cong S^1$
na mocy Stwierdzenia~\ref{prop:brzeg-sympleksu}.
Jeżeli dołączymy jedyny $2$-sympleks $V$,
otrzymamy kompleks $K_D$ z realizacją
$|K_D|=\Delta^2\cong D^2$.
O dodaniu wnętrza decyduje więc obecność
jednego zbioru $V$ w rodzinie sympleksów;
same trzy krawędzie już tworzą zamknięty obieg.
\end{przyklad}

\begin{uwaga}[Nie każdy podział na trójkąty
jest kompleksem symplicjalnym]
\label{uwaga:delta-motywacja}
Kwadrat z utożsamionymi przeciwległymi
bokami przedstawia torus. Gdy podzielimy
go jedną przekątną, po sklejeniu wszystkie
cztery narożniki stają się jednym
wierzchołkiem, a każda z dwóch
trójkątnych części nadal ma sens.
Nie jest to \emph{kompleks symplicjalny}
z jednym wierzchołkiem: trójkąt
potrzebuje trzech różnych wierzchołków.
Wygodniejszym opisem są
$\Delta$-kompleksy.
\end{uwaga}

\begin{definicja}[$\Delta$-kompleks]
\label{def:delta-kompleks}
\emph{$\Delta$-kompleks} $X$ jest
przestrzenią Hausdorffa z rodziną ciągłych odwzorowań charakterystycznych
$\sigma_\alpha:\Delta^{n_\alpha}\to X$:
\begin{enumerate}[label=(\roman*),leftmargin=*]
\item ograniczenie $\sigma_\alpha$ do wnętrza
      jest iniektywne, a obrazy wnętrz
      różnych sympleksów są rozłączne
      i pokrywają $X$;
\item ograniczenie $\sigma_\alpha$ do każdej
      ściany jest jednym z odwzorowań
      charakterystycznych niższego wymiaru
      po standardowym afinicznym
      utożsamieniu tej ściany z sympleksem, zachowującym
      kolejność pozostałych wierzchołków;
\item $A\subseteq X$ jest domknięty dokładnie
      wtedy, gdy $\sigma_\alpha^{-1}(A)$
      jest domknięty w każdym
      $\Delta^{n_\alpha}$.
\end{enumerate}
Różne ściany tego samego sympleksu
mogą być tu sklejone, o ile zgodne
są odwzorowania na ich dalszych ścianach.
\end{definicja}

Ostatni warunek to \emph{topologia słaba}
względem odwzorowań charakterystycznych:
odwzorowanie z $X$ do innej przestrzeni
jest ciągłe dokładnie wtedy, gdy jego
złożenia z każdą $\sigma_\alpha$ są ciągłe.
Jest to topologia ilorazowa z
Definicji~\ref{def:topologia-ilorazowa}
na rozłącznej sumie wszystkich sympleksów.

\begin{przyklad}[Okrąg i torus jako $\Delta$-kompleksy]
\label{prz:delta-torus}
Okrąg ma jeden $0$-sympleks $v$
i jeden $1$-sympleks, którego oba końce
przechodzą w $v$. Jego odwzorowanie charakterystyczne
$\Delta^1\to S^1$ jest iniektywne we wnętrzu,
choć na całym odcinku utożsamia oba końce;
obie ściany zerowymiarowe ograniczają się do
tego samego odwzorowania $\Delta^0\to S^1$
o obrazie $\{v\}$.
Dla torusa
podzielmy kwadrat $[0,1]^2$
przekątną i sklejmy
$(0,t)\sim(1,t)$ oraz
$(s,0)\sim(s,1)$. Otrzymujemy
jeden wierzchołek, trzy krawędzie
(poziomą $a$, pionową $b$ i przekątną
$c$) oraz dwa $2$-sympleksy. Porządkujemy ich wierzchołki jako
$((0,0),(0,1),(1,1))$ i $((0,0),(1,0),(1,1))$.
Oba porządki nadają krawędziom poziomym kierunek w prawo,
pionowym w górę, a przekątnej kierunek od $(0,0)$ do $(1,1)$.
Odwzorowania ścian są zatem zgodne także z porządkiem wierzchołków.
Znakowa reguła z Uwagi~\ref{uwaga:orientacja-przed-homologia}
pokazuje więcej:
oba zorientowane trójkąty mają algebraiczny brzeg
$a+b-c$. Górny trójkąt ma jednak orientację
przeciwną do dolnego względem orientacji kwadratu,
więc różnica dolnego i górnego trójkąta ma brzeg zero.
To konkretny przykład znoszenia się brzegów
w łańcuchu opisującym zamkniętą powierzchnię.
Iloraz jest homeomorficzny z
$(\R/\Z)\times(\R/\Z)$:
odwzorowanie $(s,t)\mapsto([s],[t])$
identyfikuje dokładnie wskazane
punkty brzegowe, jest ciągłe
i przechodzi na bijekcję
między przestrzenią zwartą
a przestrzenią Hausdorffa.
\end{przyklad}

\begin{figure}[!htbp]
\centering
\begin{tikzpicture}[>=Stealth,line cap=round,line join=round,
 x=1cm,y=1cm]
 \coordinate (a) at (0,0);
 \coordinate (b) at (4.4,0);
 \coordinate (c) at (4.4,4.4);
 \coordinate (d) at (0,4.4);
 \fill[manifoldblue!8] (a)--(b)--(c)--cycle;
 \fill[manifoldgreen!10] (a)--(d)--(c)--cycle;
 \draw[manifoldblue,very thick] (a)--(b)--(c)--(d)--cycle;
 \draw[manifoldgold!90!black,thick] (a)--(c);
 \draw[manifoldgold!90!black,thick,-{Stealth[length=2.7mm]}]
  (1.8,1.8)--(2.7,2.7);
 \draw[manifoldred,very thick,-{Stealth[length=2.7mm]}]
  (1.3,0)--(2.6,0);
 \draw[manifoldred,very thick,-{Stealth[length=2.7mm]}]
  (1.3,4.4)--(2.6,4.4);
 \draw[manifoldgreen!75!black,very thick,-{Stealth[length=2.7mm]}]
  (0,1.3)--(0,2.6);
 \draw[manifoldgreen!75!black,very thick,-{Stealth[length=2.7mm]}]
  (4.4,1.3)--(4.4,2.6);
 \node[manifoldred,below=2pt] at (3.2,0) {$a$};
 \node[manifoldgreen!65!black,right=2pt] at (4.4,3.2) {$b$};
 \node[manifoldgold!65!black,above left=2pt]
  at (2.6,2.6) {$c$};
 \foreach \p in {a,b,c,d}
  \fill[manifoldblue] (\p) circle(2.3pt);
 \node[font=\small] at (1.42,2.48) {$\Delta^2_1$};
 \node[font=\small] at (3.17,1.74) {$\Delta^2_2$};
\end{tikzpicture}
\caption{Torus z dwóch sympleksów: boki o tych samych
kolorach i kierunkach sklejamy ze sobą. Przekątna
pozostaje trzecią krawędzią. Wszystkie cztery narożniki
reprezentują ten sam wierzchołek torusa.}
\label{fig:delta-torus}
\end{figure}

\section{Grafy i drzewa}
\label{sec:grafy-drzewa}
\index{graf}
\index{drzewo!grafowe}

Graf $\Gamma=(V,E)$ składa się ze zbioru \emph{wierzchołków}
$V$ i zbioru \emph{krawędzi} $E$. Każdej krawędzi przypisujemy
dwa końce w $V$, nie wybierając jeszcze kierunku przejścia.
Końce mogą być tym samym wierzchołkiem (\emph{pętla}), a dwie
różne krawędzie mogą mieć te same końce (\emph{krawędzie
równoległe}). Dzięki temu definicja obejmuje grafy uchwytów,
grafy dualne triangulacji i jednowymiarowe szkielety CW.
\emph{Drogą krawędziową} od $v_0$ do $v_k$ jest ciąg
$v_0,e_1,v_1,\ldots,e_k,v_k$, w którym $e_i$ łączy
$v_{i-1}$ z $v_i$; jej długość wynosi $k$.
Od drogi ciągłej z Sekcji~\ref{sec:homotopie-pi1} odróżnia ją
to, że zapisuje tylko przejścia między wierzchołkami.
Droga krawędziowa jest \emph{prosta}, gdy jej wierzchołki
są różne. \emph{Cykl} to zamknięta droga krawędziowa
o różnych krawędziach, bez powtórzeń wierzchołków poza
wspólnym początkiem i końcem. Pętla jest cyklem długości
$1$, a dwie krawędzie równoległe tworzą cykl długości $2$.
Graf jest \emph{spójny}, gdy każde dwa wierzchołki
łączy droga krawędziowa.
Jest \emph{lokalnie skończony}, gdy z każdego wierzchołka
wychodzi tylko skończenie wiele krawędzi (pętlę liczymy
dwukrotnie jako dwa końce tej samej krawędzi).
\index{graf!lokalnie skończony}

\begin{definicja}[Drzewo, las i korzeń]
\label{def:drzewo-grafowe}
\emph{Drzewo} to spójny graf bez cykli. \emph{Las} to graf
bez cykli, czyli rozłączna suma drzew. W \emph{drzewie
zakorzenionym} wyróżniamy wierzchołek zwany korzeniem;
odległość od niego, liczona liczbą krawędzi jedynej
drogi, jest poziomem wierzchołka. Wierzchołek o jednej
krawędzi nazywamy \emph{liściem}, gdy nie jest korzeniem.
Korzeń nie zmienia samego drzewa; pozwala mówić o drodze
od początku i porządkować wierzchołki według poziomów.
\end{definicja}

Na przykład droga z czterech wierzchołków jest drzewem
o trzech krawędziach. Trójkąt jest spójny, lecz jego
trzy krawędzie tworzą cykl; usunięcie jednej daje drzewo.

\begin{stwierdzenie}[Drogi, krawędzie i drzewa rozpinające]
\label{prop:drzewo-rozpinajace}
W drzewie dowolne dwa wierzchołki łączy dokładnie jedna
prosta droga krawędziowa. Skończone drzewo o $v\geq1$
wierzchołkach ma $v-1$ krawędzi; skończony las o $c$
składowych ma $|E|=|V|-c$.
Każdy skończony spójny graf ma \emph{drzewo rozpinające}
$T$, czyli podgraf będący drzewem i zawierający wszystkie
wierzchołki. Każda krawędź spoza $T$ tworzy z jedyną
drogą w $T$ między jej końcami dokładnie jeden cykl.
Liczba takich dodatkowych krawędzi wynosi $|E|-|V|+1$.
\end{stwierdzenie}
\begin{proof}
Spójność daje drogę między dwoma wierzchołkami;
usuwając z niej odcinki między powtórzonymi
wierzchołkami, otrzymujemy drogę prostą.
Gdyby istniały dwie różne drogi proste o tych samych
końcach, po pierwszym rozdzieleniu musiałyby po raz
pierwszy ponownie spotkać się najpóźniej w końcu.
Ich fragmenty między tymi spotkaniami utworzyłyby
cykl, sprzecznie z definicją drzewa.

Dla $v=1$ liczba krawędzi wynosi $0$.
Przy $v>1$ wybierzmy najdłuższą drogę prostą.
Jej końcowy wierzchołek ma tylko krawędź do
poprzednika: krawędź do wierzchołka już na tej
drodze dałaby cykl, a krawędź do nowego
wierzchołka przedłużyłaby drogę.
Po usunięciu tego liścia i jego krawędzi graf
pozostaje spójny, bo droga między pozostałymi
wierzchołkami nie mogła przechodzić przez liść,
i nadal nie ma cykli. Indukcja daje
$(v-1)-1$ krawędzi w mniejszym drzewie,
a po przywróceniu liścia dokładnie $v-1$.
Sumując ten wzór po $c$ składowych lasu,
dostajemy $|E|=|V|-c$.

Aby zbudować $T$, zacznijmy od jednego wierzchołka.
Dopóki nie objęliśmy wszystkich wierzchołków,
spójność grafu zapewnia krawędź z już objętego
zbioru do wierzchołka poza nim. Dodajmy tę
krawędź i nowy wierzchołek. Każdy krok zachowuje
spójność i nie tworzy cyklu, więc po skończenie
wielu krokach otrzymujemy drzewo rozpinające.
Jeśli $e\notin T$ ma końce $u,w$, jedyna droga
w $T$ od $u$ do $w$ wraz z $e$ zamyka cykl.
Jest to jedyny cykl w $T\cup\{e\}$, bo każdy
cykl musi używać $e$. Ponieważ $T$ ma $|V|-1$
krawędzi, poza nim pozostaje $|E|-|V|+1$ krawędzi.
\end{proof}

\begin{przyklad}[Kwadrat z przekątną]
Na Rysunku~\ref{fig:graf-drzewo-cykle} graf ma cztery
wierzchołki i pięć krawędzi. Trzy niebieskie krawędzie
tworzą drzewo rozpinające. Czerwona krawędź $v_4v_1$
zamyka cykl kwadratu, a czerwona przekątna $v_1v_3$
zamyka cykl trójkąta z drogą $v_1v_2v_3$ w drzewie.
Liczba krawędzi spoza drzewa to $5-4+1=2$.
\end{przyklad}

\begin{figure}[!htbp]
\centering
\begin{tikzpicture}[line cap=round,every node/.style={font=\small}]
 \coordinate (v1) at (0,0);
 \coordinate (v2) at (3.4,0);
 \coordinate (v3) at (3.4,2.2);
 \coordinate (v4) at (0,2.2);
 \draw[manifoldblue,very thick] (v1)--(v2)--(v3)--(v4);
 \draw[manifoldred!85!black,very thick,densely dashed]
  (v4)--(v1) (v1)--(v3);
 \foreach \v in {v1,v2,v3,v4}
  \fill[manifoldgold!85!black] (\v) circle (2.5pt);
 \node[below left] at (v1) {$v_1$};
 \node[below right] at (v2) {$v_2$};
 \node[above right] at (v3) {$v_3$};
 \node[above left] at (v4) {$v_4$};
\end{tikzpicture}
\caption{Niebieskie krawędzie tworzą drzewo rozpinające $T$.
Każda z dwóch czerwonych krawędzi wraz z jedyną drogą
w $T$ między jej końcami wyznacza cykl.}
\label{fig:graf-drzewo-cykle}
\end{figure}

W drzewach nieskończonych użyjemy tylko jedyności
skończonych dróg krawędziowych; wzory na liczby
wierzchołków i krawędzi dotyczą grafów skończonych.
Graf $\Gamma$ można zrealizować topologicznie jako
$|\Gamma|$, dołączając do każdego $e\in E$ kopię
odcinka $[0,1]$ i sklejając jej końce z przypisanymi
wierzchołkami. Pętla daje wtedy okrąg z jednym
wierzchołkiem. Podział krawędzi przez dodanie punktów
wewnętrznych zmienia opis grafowy, lecz nie przestrzeń
$|\Gamma|$; aby pozbyć się pętli, wystarczy wstawić do
każdej z nich dwa nowe wierzchołki, a aby rozdzielić
krawędzie równoległe, po jednym nowym wierzchołku
we wnętrzu każdej takiej krawędzi.

Krawędzie spoza drzewa rozpinającego okażą się wolnymi
generatorami grupy podstawowej grafu.

\begin{twierdzenie}[Grupa podstawowa skończonego grafu]
\label{thm:pi1-graf}
Niech $\Gamma$ będzie skończonym spójnym grafem, $v_0$ jego
wierzchołkiem, a $T$ drzewem rozpinającym. Dla każdej krawędzi
$e\notin T$ wybierzmy kierunek od $u_e$ do $w_e$. Niech
$\gamma_e$ biegnie w $T$ od $v_0$ do $u_e$, przechodzi wzdłuż
$e$ i wraca w $T$ od $w_e$ do $v_0$. Wtedy klasy
$[\gamma_e]$ tworzą bazę wolną $\pi_1(|\Gamma|,v_0)$.
W szczególności
$\pi_1(|\Gamma|,v_0)\cong F_{|E|-|V|+1}$.
Jeśli $\Gamma$ jest drzewem, przestrzeń $|\Gamma|$ jest
ściągalna.
\end{twierdzenie}
\begin{proof}
Najpierw pokażmy ostatnie zdanie. Skończone drzewo z co
najmniej dwiema krawędziami ma liść różny od $v_0$.
Jego krawędź wraz z końcowym wierzchołkiem można
stopniowo ściągnąć do drugiego końca krawędzi, nie ruszając
reszty grafu. To jest silna retrakcja deformacyjna, ponieważ
liść nie należy do żadnej innej krawędzi. Powtarzając ten
krok, ściągamy $|T|$ do $v_0$. Dla drzewa z jedną krawędzią
wykonujemy ten sam krok, a drzewo bez krawędzi jest punktem.

Ponumerujmy krawędzie spoza drzewa jako $e_1,\ldots,e_r$
i połóżmy $G_0=|T|$ oraz
$G_j=G_{j-1}\cup e_j$ dla $1\leq j\leq r$.
Niech $p_j\subseteq T$ będzie jedyną drogą od
$u_{e_j}$ do $w_{e_j}$; dla pętli jest to droga stała.
Podgraf
$p_j\cup e_j$ jest okręgiem, również wtedy, gdy $e_j$
jest pętlą albo krawędzią równoległą do krawędzi drzewa.

Zastosujmy twierdzenie van Kampena do $G_j$, najpierw
z punktem bazowym $u_{e_j}\in p_j$. Jako pierwszy
zbiór otwarty wybieramy $G_{j-1}$ wraz z krótkimi fragmentami
$e_j$ przy jej końcach. Zbiór ten deformacyjnie retrahuje się
na $G_{j-1}$. Drugi zbiór jest małym otwartym otoczeniem
$p_j\cup e_j$ w grafie $G_j$; oprócz tych krawędzi zawiera
tylko krótkie fragmenty krawędzi incydentnych do
wierzchołków $p_j$. Retraktuje się on na okrąg
$p_j\cup e_j$, a przecięcie obu zbiorów retraktuje się
na drzewo $p_j$. Oba zbiory i ich przecięcie są spójne
drogowo. Ponieważ $\pi_1(p_j)=0$, twierdzenie van Kampena
daje $\pi_1(G_j,u_{e_j})\cong
\pi_1(G_{j-1},u_{e_j})*\mathbb Z$.
Nowy czynnik jest generowany przez przejście po $e_j$
i powrót wzdłuż $p_j$. Przenieśmy punkt bazowy do $v_0$
jedyną drogą $q$ w $T$ od $v_0$ do $u_{e_j}$.
Otrzymujemy pętlę $q*e_j*\bar p_j*\bar q$.
W drzewie droga $\bar p_j*\bar q$ ma te same końce co
jedyna droga od $w_{e_j}$ do $v_0$, a obie są homotopijne
względem końców po skasowaniu odcinków przebytych tam
i z powrotem. Jest to właśnie $[\gamma_{e_j}]$.

Zaczynając od trywialnej grupy $\pi_1(G_0)$ i dodając
krawędzie kolejno, otrzymujemy grupę wolną o bazie
$[\gamma_{e_1}],\ldots,[\gamma_{e_r}]$.
Z poprzedniego stwierdzenia $r=|E|-|V|+1$.
\end{proof}

W grafie z Rysunku~\ref{fig:graf-drzewo-cykle} dwie czerwone
krawędzie odpowiadają więc dwóm wolnym generatorom:
$\pi_1(|\Gamma|)\cong F_2$. Samo drzewo ma grupę
podstawową trywialną. Ta sama konstrukcja działa dla
każdej składowej grafu osobno. Później, gdy wprowadzimy
homologie, liczba $|E|-|V|+1$ będzie także pierwszą liczbą
Bettiego spójnego grafu.

Każdą ciągłą pętlę w skończonym grafie można homotopijnie
zastąpić skończoną drogą krawędziową. Po podziale krawędzi
na połówki otwarte gwiazdy wierzchołków tworzą pokrycie
grafu zbiorami ściągalnymi. Zwartość przedziału parametru
pozwala podzielić pętlę na skończenie wiele odcinków, z
których każdy leży w jednej gwieździe. W każdym odcinku
zastępujemy fragment pętli drogą przez jej środek; homotopia
pozostaje w tej samej gwieździe. Po połączeniu fragmentów
i usunięciu powrotów otrzymujemy drogę krawędziową.
Jeśli punkt bazowy jest wierzchołkiem, pierwszy i ostatni
fragment dobieramy tak, by homotopia zachowała ten punkt.

\section{CW-kompleksy i dołączanie komórek}
\label{sec:cw-kompleksy}

Trójkąty są wygodne do obliczeń, lecz
potrafią wymagać wielu wierzchołków.
Komórka pozwala dołączyć naraz cały
otwarty dysk, przy czym jego brzeg
może biec po starszej części przestrzeni
w sposób nietrywialny.

\begin{definicja}[Dołączenie komórki]
\index{dolaczenie komorki@Dołączenie komórki}
\label{def:dolaczenie-komorki}
Niech $Y$ będzie przestrzenią i
$\varphi:S^{n-1}=\partial D^n\to Y$
odwzorowaniem ciągłym, $n\geq1$.
\emph{Dołączeniem $n$-komórki} jest
przestrzeń ilorazowa
\[
 Y\cup_\varphi D^n
 :=(Y\sqcup D^n)/(z\sim\varphi(z)
 \text{ dla }z\in S^{n-1}).
\]
Obraz wnętrza dysku oznaczamy $e^n$.
Odwzorowanie z dysku do ilorazu jest
\emph{odwzorowaniem charakterystycznym}
tej komórki; $\varphi$ nazywamy
\emph{odwzorowaniem dołączającym}.
W wymiarze $0$ po prostu dołączamy
izolowany punkt, gdyż
$S^{-1}=\varnothing$.
\end{definicja}

Naturalna inkluzja $Y$ w przestrzeń dołączoną jest
homeomorfizmem na domknięty obraz. Dla każdego domkniętego $F\subseteq Y$
przeciwobraz jego obrazu w $Y\sqcup D^n$ jest równy
$F\sqcup\varphi^{-1}(F)$ i jest domknięty; dotyczy to również $F=Y$.
Ponieważ na punktach $Y$ nie wprowadzamy nowych utożsamień,
otrzymujemy to stwierdzenie. Wnętrze $D^n$ jest natomiast
otwartym zbiorem nasyconym, bez żadnych utożsamień, więc przechodzi
homeomorficznie na otwartą komórkę $e^n$.

Już dla $n=1$ sposób dołączenia zmienia wynik. Jeśli
$Y$ składa się z dwóch punktów i dwa końce $D^1$ trafiają
do różnych punktów, powstaje odcinek. Gdy $Y$ jest jednym
punktem i oba końce trafiają do niego, powstaje okrąg.
W obu przypadkach wnętrze dysku jest tą samą otwartą
$1$-komórką; całą różnicę zapisuje odwzorowanie
dołączające na $S^0$.

Topologia ilorazowa z
Definicji~\ref{def:topologia-ilorazowa}
ma tu konkretne znaczenie. Jeśli
$F:Y\cup_\varphi D^n\to Z$
ma być ciągła, wystarczy podać
ciągłe $F_Y:Y\to Z$ i
$F_D:D^n\to Z$ zgodne na brzegu:
$F_D(z)=F_Y(\varphi(z))$.
Istotnie, takie odwzorowania dają ciągłe
odwzorowanie na rozłącznej sumie, stałe
na każdej klasie sklejenia;
uniwersalna własność ilorazu
przenosi je na iloraz.
To jest praktyczny sposób
konstruowania odwzorowań na przestrzeniach
sklejonych.

\begin{definicja}[CW-kompleks]
\index{cw-kompleks@CW-kompleks}
\label{def:cw-kompleks}
\emph{CW-kompleks} $X$ jest przestrzenią Hausdorffa budowaną
ze szkieletów $X^0\subseteq X^1
\subseteq X^2\subseteq\cdots$:
$X^0$ jest dyskretnym zbiorem punktów,
a $X^n$ powstaje z $X^{n-1}$
przez dołączenie rodziny $n$-dysków
odwzorowaniami $S^{n-1}\to X^{n-1}$.
Przestrzeń $X=\bigcup_n X^n$
ma \emph{topologię słabą}:
$A\subseteq X$ jest domknięty,
dokładnie wtedy, gdy przeciwobraz $A$ pod każdym
odwzorowaniem charakterystycznym $D^n\to X$
jest domknięty. Wymagamy ponadto
\emph{skończoności domknięć komórek}:
domknięcie każdej komórki
spotyka tylko skończenie wiele
komórek. Litery CW przypominają
te dwa warunki
(\emph{closure-finite}, \emph{weak topology}).
\emph{Podkompleks CW} to suma komórek
domknięta w $X$; jeśli zawiera komórkę,
zawiera także całe jej domknięcie.
\end{definicja}

\begin{figure}[!htbp]
\centering
\begin{tikzpicture}[>=Stealth,line cap=round,line join=round]
 \node[font=\bfseries] at (0,3.02) {$X^0$};
 \node[font=\bfseries] at (4.33,3.02) {$X^1$};
 \node[font=\bfseries] at (8.68,3.02) {$X^2$};
 % 0-komórka.
 \fill[manifoldred] (0,1.39) circle(3.7pt);
 \node[below=7pt] at (0,1.39) {$e^0=\{v\}$};
 % Dołączenie 1-komórki: oba końce trafiają do v.
 \draw[manifoldblue,line width=2.5pt]
  (4.33,1.45) circle(1.03);
 \fill[manifoldred] (4.33,.42) circle(3.7pt);
 \node[manifoldblue!85!black] at (4.33,1.45) {$e^1$};
 \draw[->,manifoldblue!75!black,thin]
  (4.62,1.74)--({4.33+1.03*cos(45)},{1.45+1.03*sin(45)});
 \node[below=7pt] at (4.33,.42) {$v$};
 \draw[->,manifoldgold!80!black,thick]
  (1.22,1.54)--(3.08,1.54)
  node[midway,above=5pt,font=\small] {dołącz $e^1$};
 % Dołączenie 2-komórki: brzeg dysku biegnie po okręgu.
 \fill[manifoldgreen!18] (8.68,1.45) circle(1.03);
 \draw[manifoldblue,line width=2.5pt]
  (8.68,1.45) circle(1.03);
 \fill[manifoldred] (8.68,.42) circle(3.7pt);
 \node[manifoldgreen!65!black] at (8.68,1.45) {$e^2$};
 \node[below=7pt] at (8.68,.42) {$v$};
 \draw[->,manifoldgold!80!black,thick]
  (5.58,1.54)--(7.35,1.54)
  node[midway,above=5pt,font=\small] {dołącz $e^2$};
\end{tikzpicture}
\caption{Przykłady kolejnych wymiarów CW: punkt
$X^0=\{v\}$, okrąg $X^1\cong S^1$ otrzymany
przez sklejenie obu końców $1$-komórki z $v$,
oraz dysk $X^2\cong D^2$ po dołączeniu
$2$-komórki wzdłuż całego okręgu.
Czerwony punkt, niebieski łuk i zielone wnętrze
oznaczają odpowiednio otwarte komórki
wymiarów $0,1,2$.}
\label{fig:cw-wymiary}
\end{figure}

W kompleksie skończonym warunki
topologiczne łatwo sprawdzić
bez osobnej teorii granic:
topologia słaba jest ilorazową
topologią sklejania skończenie
wielu dysków, a domknięcie
każdej komórki spotyka
tylko skończenie wiele innych.
Przy nieskończenie wielu komórkach
topologia słaba jest konieczna,
aby ciągłość dało się sprawdzać
na wszystkich domkniętych komórkach.

Na przykład rozważmy przeliczalną sumę okręgów z jednym
wspólnym wierzchołkiem, zbudowaną jako CW-kompleks z
jednej $0$-komórki i przeliczalnie wielu $1$-komórek.
Wybierzmy jeden punkt $p_j$ we wnętrzu każdej komórki.
Zbiór $\{p_j:j\geq1\}$ jest domknięty w topologii słabej:
jego przeciwobraz w każdym dysku charakterystycznym jest
jednym punktem, więc jest domknięty. W szczególności
$p_j$ nie zbiegają do wspólnego wierzchołka. Gdy natomiast
rysujemy okręgi w płaszczyźnie z promieniami malejącymi do
zera, punkty położone naprzeciw wspólnego punktu zbiegają
do niego. Ten rysunek ma inną topologię przy wierzchołku;
warunek słabej topologii rozstrzyga tę różnicę.

\begin{stwierdzenie}[Sympleksy dają komórki]
\label{prop:sympleksy-cw}
Skończony kompleks symplicjalny i
skończony $\Delta$-kompleks są
CW-kompleksami: ich otwarte komórki
to obrazy wnętrz sympleksów.
\end{stwierdzenie}
\begin{proof}
Indukujemy po wymiarze. Szkielet
zerowy składa się z wierzchołków.
Wnętrze każdego $n$-sympleksu
jest rozłączne z obrazami innych
wnętrz, a brzeg jest sumą ścian
zawartych w poprzednim szkielecie.
Stwierdzenie~\ref{prop:brzeg-sympleksu}
pozwala traktować sympleks
jako dysk $D^n$, którego
odwzorowanie na brzegu jest odwzorowaniem dołączającym.
Warunek o topologii jest w
Definicji~\ref{def:delta-kompleks}
podany wprost, a dla skończonej
realizacji symplicjalnej wynika
z topologii skończonej sumy zwartych sympleksów:
ciągła surjekcja z tej sumy na $|K|\subseteq\R^V$ jest domknięta,
gdyż obrazy zbiorów domkniętych są zwarte w przestrzeni Hausdorffa.
Jest więc odwzorowaniem ilorazowym, co daje właśnie wymaganą topologię.
Skończona liczba sympleksów
zapewnia skończoność domknięć komórek.
\end{proof}

\begin{przyklad}[Minimalne struktury komórkowe]
\label{prz:cw-sfera-torus}
\begin{enumerate}[label=(\arabic*),leftmargin=*]
\item Dla $n\geq1$ sfera $S^n$
      ma jedną $0$-komórkę i jedną
      $n$-komórkę:
      $D^n/S^{n-1}\cong S^n$.
      Jawny homeomorfizm opisujemy
      poniżej.
      Dla $S^0$ używamy dwóch
      $0$-komórek.
\item Torus ma jedną $0$-komórkę,
      dwie $1$-komórki $a,b$ oraz
      jedną $2$-komórkę. Obchodząc
      brzeg kwadratu z
      Rysunku~\ref{fig:delta-torus},
      otrzymujemy słowo
      $aba^{-1}b^{-1}$ w
      $S^1\vee S^1$.
      Dwa trójkąty realizacji
      $\Delta$-kompleksowej zostały
      zastąpione jednym dyskiem,
      którego odwzorowanie dołączające
      zapamiętuje sklejanie.
\end{enumerate}
\end{przyklad}

\begin{przyklad}[Sfera $S^2$ bez $1$-komórek]
\label{prz:cw-s2}
Weźmy $X^0=\{v\}$ i nie dołączajmy
żadnej $1$-komórki, więc $X^1=X^0$.
Niech odwzorowanie dołączające
$\varphi:\partial D^2=S^1\to\{v\}$
będzie stałe. Wtedy
\[
 X^2=X^1\cup_\varphi D^2
 \cong D^2/\partial D^2\cong S^2.
\]
Cały brzeg dysku staje się \emph{jednym}
punktem $v$; obraz jego wnętrza
jest jedną otwartą $2$-komórką
$e^2=S^2\setminus\{v\}\cong\R^2$.
Rysunek~\ref{fig:cw-s2} pokazuje to
sklejenie. W szczególności
$X^1$ nie jest okręgiem:
okrąg po lewej stronie rysunku
to brzeg \emph{dołączanego dysku},
a nie komórka przestrzeni wynikowej.
\end{przyklad}

\begin{figure}[!htbp]
\centering
\begin{tikzpicture}[>=Stealth,line cap=round,line join=round]
 % Dysk przed sklejeniem: czerwony brzeg zapada się do v.
 \fill[manifoldblue!10] (0,1.85) circle(1.33);
 \draw[manifoldred,line width=2.4pt] (0,1.85) circle(1.33);
 \node[font=\bfseries] at (0,3.58) {$D^2$};
 \node[manifoldblue!80!black] at (0,1.85) {$\operatorname{int}D^2$};
 \node[manifoldred!85!black,below=7pt]
  at (0,.52) {$\partial D^2=S^1$};
 % Sklejenie brzegu do jednego punktu.
 \draw[->,manifoldgold!85!black,line width=1.5pt]
  (1.73,1.85)--(4.8,1.85)
  node[midway,above=7pt] {$q$};
 \node[font=\small,manifoldgold!70!black]
  at (3.26,1.25) {sklejamy cały brzeg};
 % Sfera w stylistyce rysunku 1.4: jasny środek, tył i przód łuków.
 \shade[inner color=white,outer color=manifoldblue!43]
  (6.8,1.85) circle(1.43);
 \draw[manifoldblue!75!black,very thick]
  (6.8,1.85) circle(1.43);
 \draw[manifoldblue!50,densely dashed,thin]
  (8.23,1.85) arc[start angle=0,end angle=180,x radius=1.43,y radius=.34];
 \draw[manifoldblue!70!black,thin]
  (5.37,1.85) arc[start angle=180,end angle=360,x radius=1.43,y radius=.34];
 \draw[manifoldblue!38,densely dashed,thin]
  (6.8,3.28) .. controls (6.35,2.76) and (6.35,.94) .. (6.8,.42);
 \draw[manifoldblue!65!black,thin]
  (6.8,3.28) .. controls (7.25,2.76) and (7.25,.94) .. (6.8,.42);
 \fill[manifoldred] (6.8,.42) circle(3.8pt);
 \node[font=\bfseries] at (6.8,3.58) {$X^2\cong S^2$};
 \node[manifoldblue!85!black] at (6.8,2.44) {$e^2$};
 \node[manifoldred!85!black,below=7pt] at (6.8,.42) {$v$};
\end{tikzpicture}
\caption{Struktura CW sfery: $D^2$ dołączamy do
punktu $v$ odwzorowaniem stałym na całym brzegu;
$q:D^2\to D^2/\partial D^2$ jest odwzorowaniem ilorazowym.
W powstałej sferze czerwony punkt to
jedyne $e^0$, a cała pozostała powierzchnia
to jedna otwarta komórka $e^2$.
Nie występują komórki wymiaru $1$.}
\label{fig:cw-s2}
\end{figure}

\begin{uwaga}[Uwaga o ilorazie dysku]
W punkcie (1) nie potrzebujemy
identyfikować półkuli z ilorazem
za pomocą niejawnego rysunku.
Istnieje jawne ciągłe odwzorowanie
$D^n\to S^n\subseteq\R^{n+1}$,
\[
 x\longmapsto
 \bigl(\sin(\pi\|x\|)\,x/\|x\|,
       \cos(\pi\|x\|)\bigr)
 \quad(x\ne0),
\]
z wartością $(0,\ldots,0,1)$
w $x=0$. Cały brzeg $\|x\|=1$
przechodzi w biegun południowy,
a na wnętrzu odwzorowanie jest iniektywne
i pokrywa sferę bez tego bieguna.
Przechodzi więc na ciągłą bijekcję
$D^n/S^{n-1}\to S^n$, która jest
homeomorfizmem przez zwartość
dziedziny i własność Hausdorffa
sfery.
\end{uwaga}

\begin{twierdzenie}[Co robi dołączenie komórki z $\pi_1$?]
\label{thm:pi1-dolaczenie-komorki}
Niech $Y$ będzie łukowo spójna,
a $X=Y\cup_\varphi D^n$, z punktem bazowym $y_0\in Y$.
Jeśli $n=2$, wybierzmy $z_0\in S^1$, drogę
$\lambda$ w $Y$ od $y_0$ do $\varphi(z_0)$ oraz
dodatnio zorientowaną pętlę $\ell$ w $S^1$ opartą w $z_0$.
Wtedy $\pi_1(X,y_0)$ jest ilorazem $\pi_1(Y,y_0)$
przez domknięcie normalne elementu
$[\lambda*(\varphi\circ\ell)*\bar\lambda]$.
Jeśli $n\geq3$, inkluzja
$Y\hookrightarrow X$ wyznacza izomorfizm na $\pi_1$.
\end{twierdzenie}
\begin{proof}
W dysku wybierzmy promień $r=\|x\|$.
Otwarty w $X$ zbiór $U$ złożony
z $Y$ i pasa $r>1/3$ deformuje
się na $Y$: przesuwamy punkty
pasa promieniowo do $r=1$,
a punkty $Y$ zostawiamy. Na części dyskowej homotopia ma wzór
$ru\mapsto((1-t)r+t)u$, gdzie $u\in S^{n-1}$.
Wyjaśnijmy ciągłość przy sklejanym $Y$. Dla otwartego otoczenia
$O$ punktu $y\in Y$ zmniejszamy jego część dyskową do punktów $ru$,
dla których $\varphi(u)\in O\cap Y$ i cały odcinek
$\{\rho u:r\leq\rho\leq1\}$ leży w przeciwobrazie $O$.
Ta część jest otwarta: zawieranie zwartego odcinka w zbiorze
otwartym zachowuje się przy małej zmianie $r,u$.
Wraz z $O\cap Y$ tworzy otwarte otoczenie $y$ w ilorazie,
którego obraz przez homotopię pozostaje w $O$ dla wszystkich $t$.
Dowodzi to ciągłości przy $Y$; na wnętrzu dysku wynika ona z podanego wzoru.
Drugi otwarty zbiór $V$ jest
obrazem $r<2/3$; jest to
ściągalne wnętrze dysku.
Przecięcie $U\cap V$ to pierścień
$1/3<r<2/3$, homotopijnie
równoważny $S^{n-1}$.
Jest łukowo spójne dla $n\geq2$.
Otwartość obu zbiorów sprawdzamy przed sklejeniem:
przeciwobrazy $U$ to całe $Y$ i otwarty w $D^n$ pas $r>1/3$,
a przeciwobraz $V$ to $r<2/3$, bez punktów $Y$.

Twierdzenie~\ref{thm:van-kampen}
z $G_V=1$ mówi, że $\pi_1(X)$
jest ilorazem $\pi_1(U)\cong\pi_1(Y)$
przez domknięcie normalne obrazu
$\pi_1(U\cap V)$. Przy $n=2$
generator $\pi_1(S^1)\cong\Z$
po przejściu przez retrakcję
$U\to Y$ staje się pętlą
$\lambda*(\varphi\circ\ell)*\bar\lambda$.
Inna droga $\lambda$ daje element sprzężony, więc
to samo domknięcie normalne.
Przy $n\geq3$ mamy
$\pi_1(S^{n-1})=0$ z
Przykładu~\ref{prz:pi1-sfera},
więc nie pojawia się żadna
relacja ani nowy generator.
\end{proof}

\begin{wniosek}[Prezentacja grupy podstawowej kompleksu CW]
\label{wn:pi1-cw-prezentacja}
Niech $X$ będzie spójnym skończonym CW-kompleksem,
$v_0\in X^0$, a $T$ drzewem rozpinającym grafu $X^1$.
Wybierzmy kierunek każdej $1$-komórki spoza $T$.
Grupa $\pi_1(X,v_0)$ ma prezentację z jednym generatorem
na każdą taką komórkę i jedną relacją na każdą
$2$-komórkę. Brzeg jej dysku zastępujemy homotopijną
zamkniętą drogą krawędziową; relację odczytujemy z tej
drogi, zapisując krawędzie spoza $T$ ze znakami zależnymi
od kierunku i pomijając krawędzie $T$. Komórki wymiarów
co najmniej $3$ nie zmieniają tej prezentacji.
\end{wniosek}
\begin{proof}
Jednowymiarowy szkielet jest spójnym grafem: komórka
wymiaru co najmniej $2$ ma spójny brzeg, więc jej
dołączenie nie łączy różnych składowych wcześniejszego
szkieletu. Z Twierdzenia~\ref{thm:pi1-graf} grupa
$\pi_1(X^1,v_0)$ jest wolna na krawędziach spoza $T$.
Dołączajmy teraz kolejno skończenie wiele $2$-komórek.
Twierdzenie~\ref{thm:pi1-dolaczenie-komorki} dodaje dla
każdej z nich do jądra dokładnie domknięcie normalne
klasy jej brzegu. Po zastąpieniu brzegu homotopijną
pętlą krawędziową, jak wyżej, drogę od $v_0$ do jej
początkowego wierzchołka wybieramy w $T$. Zapis brzegu
w bazie z twierdzenia o grafie jest wtedy opisanym słowem;
zmiana drogi lub punktu startowego sprzęga to słowo i
nie zmienia domknięcia normalnego. Ponowne zastosowanie
tego twierdzenia do komórek wyższych wymiarów
nie zmienia grupy podstawowej.
\end{proof}

\begin{przyklad}[Torus i powierzchnia rzutowa]
\label{prz:pi1-torus-rp2}
Dla torusa część jednowymiarowa
jest $S^1\vee S^1$, z grupą
$F(a,b)$ z Przykładu~\ref{prz:pi1-dwa-okregi}.
Dołączona $2$-komórka wymusza
$aba^{-1}b^{-1}=1$, czyli
\[
 \pi_1(T^2)\cong
 \langle a,b\mid ab=ba\rangle
 \cong\Z^2.
\]
Rachunek tej prezentacji przeprowadziliśmy w
Przykładzie~\ref{prz:prezentacja-torusa-5}; dołączenie komórki
uzasadnia teraz topologicznie pochodzenie relacji.
Drugim przykładem jest
$\R P^2$: dysk będący górną
półkulą $S^2$ ma na brzegu
antypody utożsamione, więc
$\R P^2$ ma $0$-komórkę,
      $1$-komórkę i $2$-komórkę
      dołączoną do $S^1\cong\R P^1$
      odwzorowaniem odpowiadającym
      $z\mapsto z^2$ (liczba owinięć wynosi $2$).
      Istotnie, $S^1/(z\sim-z)\to S^1$, $[z]\mapsto z^2$,
      jest homeomorfizmem jako ciągła bijekcja ze zwartej
      przestrzeni do przestrzeni Hausdorffa; odwzorowanie dołączające
      przechodzi przez ten iloraz. Pełny obieg brzegu daje zatem dwa obiegi
      $1$-szkieletu, co wyjaśnia relację $a^2=1$.
      Zatem
\[
 \pi_1(\R P^2)
 \cong\langle a\mid a^2=1\rangle
 \cong\Z/2\Z.
\]
W obu przykładach relacja grupowa pochodzi bezpośrednio
z przebiegu brzegu dołączanej $2$-komórki.
\end{przyklad}

\begin{uwaga}[Granica obecnego rozdziału]
Sympleksy i komórki dają
konkretne przestrzenie, a ich
odwzorowania dołączające kodują sposób
sklejenia. W następnym rozdziale
zorientowanym sympleksom przypiszemy
generatory łańcuchów, a ich ścianom
operator brzegu $\partial$.
Sprawdzimy $\partial^2=0$,
porównamy homologię symplicjalną,
singularną i komórkową oraz
wykorzystamy długie ciągi
z Rozdziału~\ref{ch:algebra-homologiczna}.
Wyższe grupy homotopii i twierdzenie
Hurewicza wymagają już tak
zbudowanej homologii przestrzeni.
\end{uwaga}

\par\smallskip
{\footnotesize\noindent\emph{Dalsza lektura:}
Pełne rozwinięcie teorii homotopii,
nakryć, kompleksów CW i
$\Delta$-kompleksów podaje
Allen Hatcher, \emph{Algebraic Topology},
rozdziały 0 i 1 oraz początek rozdziału 2,
\url{https://pi.math.cornell.edu/~hatcher/AT/AT.pdf}.
W tym rozdziale wyprowadziliśmy
podstawowe własności, które
wykorzystamy do budowania homologii.\par}

\clearpage
